\begin{dcomment}{exposition}{p1:On est ramené aussitôt au cas où}{The presheaf represented by a change of index}
Grothendieck begins again with ordinary limits because their elementary descent argument already contains the object which the cohomological theory must derive. We keep his variance throughout: a presheaf on $C$ is a functor $C^{\mathrm{op}}\to\mathrm{Set}$, and $c/\alpha$, also written $c\backslash C'$, has objects $(a,f:c\to\alpha a)$. The arrow in the first proposition is printed in the opposite direction; its intended direction is fixed both by this notation and by the proof on manuscript page~6. This is an edition slip, not a change of convention.

Put $K_\alpha(c)=\pi_0(c/\alpha)$. Precomposition with $c_0\to c_1$ defines $K_\alpha(c_1)\to K_\alpha(c_0)$, so $K_\alpha$ is a presheaf. There is a natural bijection
\[
 \lim_{C'^{\mathrm{op}}}\alpha^*F
 \;\cong\;\operatorname{Nat}(K_\alpha,F).
\]
Indeed a compatible family $(s_a)$ determines the map $[a,f]\mapsto F(f)(s_a)$. A morphism $(a,f)\to(b,g)$ in $c/\alpha$ gives $g=\alpha(h)f$, and compatibility gives $s_a=F(\alpha h)(s_b)$; hence the value is constant along every zigzag. Naturality in $c$ follows by composition. Conversely, from $t:K_\alpha\to F$, put $s_a=t_{\alpha a}[a,1_{\alpha a}]$. For $h:a\to b$, the arrow $(a,1)\to(b,\alpha h)$ and naturality of $t$ prove compatibility. These constructions are inverse.

Equivalently, $K_\alpha=\alpha_!1_{C'}$. Thus the restriction map on limits is induced by the augmentation $K_\alpha\to1_C$. Yoneda, applied in the presheaf category, proves the exact criterion:
\[
 \text{restriction preserves every presheaf limit}
 \quad\Longleftrightarrow\quad
 \pi_0(c/\alpha)=\{*\}\quad(c\in C).
\]
The nonempty, connected comma categories are what matter. A common predecessor for each pair of objects is a stronger sufficient condition; it has no equalization requirement. For coefficients in a locally small category $\mathcal M$, apply $\operatorname{Hom}_{\mathcal M}(M,-)$ for every $M$. These functors preserve existing limits and jointly reflect isomorphisms, which justifies the reduction to sets without imposing extra completeness on $\mathcal M$.

In current covariant terminology, $\alpha$ is a final functor, while $\alpha^{\mathrm{op}}$ is initial; the latter is the functor acting on the diagrams whose limits are written here. SGA~4 uses the cofinal terminology followed in the manuscript \citep[Expos\'e~I, \S8.1]{SGA4I1972}. Proposition~3.2 below returns to precisely this criterion. We shall use the representing presheaf again there, rather than repeat the elementwise argument.
\end{dcomment}

\begin{dcomment}{exposition}{p3}{Why the relative problem has two endpoints}
The passage from global sections to a direct image introduces a second endpoint. For $f:X\to S$ and a presheaf $F$, the pointwise right Kan extension formula is
\[
 (f_*F)(s)=\lim_{(f/s)^{\mathrm{op}}}F(x),
 \qquad (f/s)=\{(x,f(x)\to s)\}.
\]
This is the ordinary comma-category formula for a Kan extension \citep[Theorem~6.2.1 and Corollary~6.2.9]{Riehl2016}. If $f'=f\alpha$, the comparison at $s$ is restriction along $\alpha/s$. Its comma category at $(x_0,u_0)$ consists exactly of factorizations
\[
 f(x_0)\xrightarrow{\,f(\xi)\,}f(\alpha a)
 \xrightarrow{\,u\,}s,
 \qquad u f(\xi)=u_0.
\]
Connectedness and nonemptiness of these categories therefore suffice. The weak cofiltering condition written here also suffices, but the phrase ``ce qui signifie aussi'' on manuscript page~1 should not turn that sufficient condition into a necessary one.

When $S=X$ and $f=1_X$, the test concerns factorizations of each arrow $x_0\to x_1$ through $\alpha a$. It is stronger than the test involving only $x_0/\alpha$. This is the same distinction already isolated in the commentary on Chapter~I, manuscript p.~49 (``The factorization category controls base change''); Proposition~4.1 will identify its ordinary categorical form. Keeping the two endpoints visible prevents an erroneous inference from invariance of global sections to full faithfulness of restriction.
\end{dcomment}

\begin{dcomment}{later}{p5:donc c'est OK pour $F$ constant.}{Cohomological cofinality and the derived weight}
Grothendieck's question about abelian coefficients has an affirmative answer. Its proof can follow the argument he has just discovered, provided that the presheaf $K_\alpha$ is replaced by its chain-level predecessor. We give the argument for a commutative ring $k$; $k=\mathbb Z$ answers the stated question. Write $D(C;k)$ for the derived category of $k$-module presheaves and use $\mathbf R\operatorname{Hom}_{C;k}$ for the \emph{global} derived Hom, valued in $D(k)$.

Derived adjunction gives, naturally in $F$,
\[
 \mathbf R\Gamma(C',\alpha^*F)
 \simeq
 \mathbf R\operatorname{Hom}_{C;k}
       (\mathbf L\alpha_!k_{C'},F).
\]
For clarity, this can be constructed without assuming the disputed internal formula on manuscript page~10. Resolve $k_{C'}$ by the augmented simplicial presheaf of free modules whose degree $n$ is
\[
 P_n=\bigoplus_{a_0\to\cdots\to a_n} k\operatorname{Hom}_{C'}(-,a_0).
\]
At $a$, its augmented chain complex is the nerve complex of $a/C'$, contracted by its initial object $1_a$. Each $P_n$ is projective. Its normalized complex, in nonpositive cohomological degrees, is therefore a projective resolution. Since restriction is exact, $\alpha_!$ takes these projectives to projectives. Evaluating $\alpha_!P_\bullet$ at $c$ gives the free chain complex of $c/\alpha$. In particular
\[
 (\mathbf L\alpha_!k_{C'})(c)
 \simeq C_*(N(c/\alpha);k),
\]
where a chain degree $n$ is placed in cohomological degree $-n$. The usual formulation using the opposite comma category gives the same complex up to signed reversal.

If every $N(c/\alpha)$ is weakly contractible, its augmented complex is quasi-isomorphic to $k$. Consequently $\mathbf L\alpha_!k_{C'}\to k_C$ is an isomorphism. Applying derived Hom yields the desired comparison in every degree, and indeed for every complex $F$. This proof needs only $k$-acyclicity of these nerves; weak contractibility is the uniform condition which works for homotopy types as well. The two conditions must not be confused with mere connectedness. For example, $BG\to e$ is cofinal for ordinary presheaf limits whenever $G$ is a group, but for $G=\mathbb Z$ and the constant coefficient $\mathbb Z$ one has $H^1(BG;\mathbb Z)=\mathbb Z$.

For spaces, the corresponding weight is $c\mapsto B(c/\alpha)$ in $\operatorname{Fun}(C^{\mathrm{op}},\mathcal S)$. The same mapping-object adjunction proves cofinality; the general higher categorical criterion is \citet[Theorem~4.1.3.1, together with Proposition~4.1.1.8]{Lurie2009HTT}, applied with the appropriate opposites. Its corollary~4.1.3.3 is Quillen's Theorem~A. These results settle the question on page~5, without requiring the stronger unit or counit assertions in its preceding paragraph. Those assertions are conditional there on the manuscript's invocation of $\mathrm{Der}_5$; interpreted as consequences of comma-category asphericity under present derivator axioms, they are false. The endpoint calculation in the next block makes the distinction explicit, continuing the caution about the changing axioms in the commentary on Chapter~I, manuscript p.~43 (``What the localization axioms do and do not say'').
\end{dcomment}

\begin{dcomment}{examples}{p8:Si par contre on travaille avec le dérivateur non additif HOT}{Global adjunction does not identify an internal Hom}
The proposed formula is useful because its global version is right. The difficulty begins only when its two sides are compared as objects over $C$. For presheaves of spaces, the correct equivalence is
\[
 \operatorname{Map}_{\mathcal P(C')}(1,\alpha^*F)
 \simeq
 \operatorname{Map}_{\mathcal P(C)}(\alpha_!1,F),
 \qquad \mathcal P(C)=\operatorname{Fun}(C^{\mathrm{op}},\mathcal S).
\]
Since $\mathcal P(C)$ is cartesian closed, the right side is also
$\mathbf R\Gamma(C,\underline{\operatorname{Hom}}(\alpha_!1,F))$.
This proves exactly the globally sectioned formula on page~8. The analogous abelian formula on page~7 takes values in $D(k)$ when ``Hom internes'' means the $k$-enriched global Hom. A sheaf-valued internal Hom would require an additional $\mathbf R\Gamma(C,-)$.

To see what global sections forget, take $C=[1]=(0\to1)$ and $i:e\to C$ with $i(*)=1$. A presheaf is an arrow $F_1\to F_0$. Direct calculation of all three adjoints gives
\[
 i^*F=F_1,\qquad i_!A=(A\xrightarrow{1}A),
 \qquad i_*A=(A\to *).
\]
Here the displayed order is always the restriction map from value at $1$ to value at $0$. Every $c/i$ is terminal, so $i$ passes even the contractible-comma criterion; $\Gamma(C,F)=F_1$ and the comparison with $\Gamma(e,i^*F)$ is an identity. Nevertheless the unit has component $F_0\to *$ at $0$, and the counit $i_!i^*F\to F$ has component $F_1\to F_0$ there. Neither is generally invertible.

Moreover, $i_!1=1_C$, hence
\[
 \underline{\operatorname{Hom}}_C(i_!1,F)=F,
 \qquad i_*i^*F=(F_1\to *).
\]
Taking $F=(\{*\}\to\{0,1\})$ disproves the internal identification. Applying global sections to both sides gives $\{*\}$, as it should. With $j(*)=0$ instead, every $j/c$ is terminal, whereas $j_*j^*F=(F_0\xrightarrow{1}F_0)$; thus the opposite comma orientation also fails to imply the universal unit assertion. The same examples work in $D(k)$, replacing $*$ by the zero module and choosing an arrow which is not an isomorphism.

The shift argument on page~9 establishes the expected Hom groups once the derived adjunction is available. Recovering an equivalence of enriched Hom objects requires the enriched adjunction itself, or a constructed comparison and a conservativity argument. It is not supplied merely by a bijection on $\pi_0$ in the space-valued case. No such extra step is needed for the explicit bar-resolution proof above.
\end{dcomment}

\begin{dcomment}{exposition}{p9:et la formule doit résulter}{The projection formula is an extra property}
Grothendieck correctly reduces the proposed internal adjunction to a projection formula. What fails is the expectation that every change of indexing category has this property. In a closed symmetric monoidal setting, write $L=\alpha_!$, $R=\alpha^*$ and $U=\alpha_*$. The canonical map supplied by the adjunction $L\dashv R$ has direction
\[
 L(RG\otimes G')\longrightarrow G\otimes LG'.
\]
It is adjoint to $RG\otimes G'\to RG\otimes RLG'$, using the unit and the monoidal structure of $R$. The arrow printed on page~10 points in the inverse direction; such an inverse is available only after invertibility has been established.

For every $H$, the adjunctions give
\[
\begin{aligned}
 \operatorname{Hom}(H,U\underline{\operatorname{Hom}}(G',RF))
 &\cong\operatorname{Hom}(L(RH\otimes G'),F),\\
 \operatorname{Hom}(H,\underline{\operatorname{Hom}}(LG',F))
 &\cong\operatorname{Hom}(H\otimes LG',F).
\end{aligned}
\]
Thus the canonical internal comparison
$\underline{\operatorname{Hom}}(LG',F)\to
U\underline{\operatorname{Hom}}(G',RF)$
is invertible for all $F,G'$ exactly when the projection maps are invertible for all $H,G'$. This proves the reduction, and also locates the additional hypothesis precisely.

For cartesian products this property is called Frobenius reciprocity. An essential geometric morphism need not satisfy it. Moreover, Frobenius reciprocity itself is weaker than its validity in every slice, the latter being the condition appropriate to a locally connected geometric morphism. \citet[\S1]{GarnerStreicher2021} explicitly distinguish these conditions; they also recall the equivalence between Frobenius reciprocity and preservation of exponentials by the right adjoint of an adjunction between cartesian closed categories. These later distinctions explain why existence of both Kan extensions does not settle the question. They should not be replaced by an unrestricted ``projection formula for derivators''.
\end{dcomment}

\begin{dcomment}{development}{p10:On trouve par adjonction un homomorphisme canonique}{The exact obstruction for a point inclusion}
The walking-arrow example can be sharpened to measure exactly which restriction maps obstruct the projection formula. This makes the question on page~10 concrete even before a general theory of Frobenius adjunctions is introduced.

\textbf{Point-inclusion criterion.}\ours\
Let $d$ be an object of a small category $C$, and let $i:e\to C$ select $d$. For a set-valued presheaf $G$, the projection map
\[
 i_!(i^*G\times A)\longrightarrow G\times i_!A
\]
is invertible for every set $A$ if and only if $G(f):G(d)\to G(c)$ is bijective for every arrow $f:c\to d$. Consequently the formula holds for all $G,A$ if and only if every arrow with codomain $d$ is an isomorphism.

\textit{Proof.}
The left adjoint to evaluation at $d$ is $i_!A=A\times h_d$, where $h_d(c)=C(c,d)$. At $c$ the projection map is
\[
 A\times G(d)\times C(c,d)\longrightarrow
 A\times G(c)\times C(c,d),\qquad
 (a,s,f)\longmapsto(a,G(f)s,f).
\]
For $A=1$ it is a disjoint union, indexed by $f$, of the maps $G(f)$; it is bijective precisely when each of them is. This also suffices for every $A$. If every $f$ is invertible, every $G(f)$ is bijective. Conversely, suppose the formula holds for every presheaf. Taking $G=h_z$ gives a bijection $C(d,z)\to C(c,z)$ by precomposition with $f$, for every $z$. With $z=c$, obtain $g:d\to c$ such that $gf=1_c$. With $z=d$, injectivity and $(fg)f=f=1_df$ imply $fg=1_d$. Thus $f$ is invertible.\hfill$\square$

This includes the object $1$ of $[1]$, where the noninvertible incoming arrow is exactly the obstruction. It also explains why the analogous formula for a group action behaves differently: all indexing arrows are invertible. For a fixed $G$, local constancy along the arrows entering $d$ is the precise repair; no conclusion about all coefficients follows from connectedness of a comma category.

We checked the pointwise Kan-extension treatment of \citet[\S6.2]{Riehl2016} and the Frobenius discussion of \citet[\S1]{GarnerStreicher2021}, and searched for the point-evaluation formulation. We could not find this elementary criterion stated there in this form. It is recorded as our worked continuation of the calculation, with no claim of priority.
\end{dcomment}

\begin{dcomment}{context}{p10}{Why the argument passes through toposes}
The change of subject at the beginning of \S2 is part of the cofinality problem. Grothendieck wants conditions on indexing categories which survive passage to categories of coefficients and which force the relevant homotopy types to disappear. A formulation in terms of the presheaf topos can detect all objects over the category, whereas the underlying preorder remembers only whether an arrow exists.

This is a small-scale instance of the program described in \citet[\S69, especially the discussion of diagram categories on pp.~1--3 of the excerpt]{Grothendieck1983}: the derived category of coefficients is accompanied by derived categories of diagrams of every shape. The concern there comes from operations on coherent diagrams which a single triangulated category does not retain. The present calculation asks what changing the shape does to sections and their derived functors. In parallel, \citet[Chapter~II, \S1; Chapter~III, \S2]{Heller1988} axiomatizes homotopy theories and develops the computation of homotopy Kan extensions. These references explain the setting, not a historical claim that one author's argument was taken from the other.

The additional tensor language on pages~7--10 needs separate structure. Heller's Chapter~IV treats tensor and cotensor operations; the preceding projection counterexample shows why that structure cannot be inferred from the mere existence of left and right Kan extensions. We now return to a question for which the ordinary presheaf topos is already sufficient: which part of the filtering conditions is visible in the behavior of its connected objects?
\end{dcomment}

\begin{dcomment}{exposition}{p11:tout sous-faisceau d'un faisceau connexe est connexe}{Irreducibility of an object and of its support}
Grothendieck is distinguishing two ways to forget morphisms. The space $\operatorname{Esp}(X)$ remembers downward-closed sets of objects; the topos $\widehat X=\operatorname{Fun}(X^{\mathrm{op}},\mathrm{Set})$ also remembers every category of elements. The first sees a common predecessor of two objects. The second must see common predecessors compatible with prescribed arrows.

Here an object is \emph{connected} if it has no decomposition into two noninitial disjoint summands; ``$0$-connected'' adds noninitiality. Thus the assertion that every subobject of a connected object is connected allows the initial subobject. An object $A$ is irreducible when it is noninitial and any two noninitial subobjects of $A$ meet. Equivalently, every subobject is connected: a disconnected subobject gives two disjoint nonempty subobjects, while disjoint subobjects have a disconnected union. Applied to $A=1$, this gives the definition of an irreducible topos. Applied to every nonempty connected $A$, after assuming local connectedness, it gives the local definition.

For $\widehat X$, subterminal objects are precisely sieves of objects: if $x$ belongs and $y\to x$ exists, then $y$ belongs. Any two nonempty such sieves meet exactly when each pair of objects has a common predecessor. This proves
\[
 \widehat X\text{ irreducible}
 \quad\Longleftrightarrow\quad
 X\ne\varnothing\text{ and }
 (\forall x,y)(\exists z)(z\to x, z\to y).
\]
It is a condition on the associated preorder, not on equalities of arrows. A least class in that preorder need not give an initial object: uniqueness of arrows has been discarded. If the preorder quotient is artinian, a minimal class exists, and common lower bounds force it below every other class; this proves the assertion on page~13. Finite principal lower sets suffice as well. ``Locally finite'' here must have that meaning, as indicated by the finite slices $X/x$; finiteness of intervals alone would fail, for example for the ordered set $\mathbb Z$.

The list of equivalences on pages~12--13 is provisional. Its claim about components being weakly cofiltered is too weak, and its subsequent inclusion of PS~2 is too strong. The annotations and Proposition~2.1 correct these in opposite directions. The surviving characterization is PS~1 alone. We prove it below, retaining the successive formulations as part of the text's development.
\end{dcomment}

\begin{dcomment}{examples}{p12:les composantes connexes de $T$ sont irréductibles}{The generic fibre without a generic point}
The topological case supplies the intuition for the definition, but it also calls for a precise interpretation of ``generic fibre''. An irreducible space need not possess a generic point. For a sheaf $F$ on a nonempty irreducible space $T$, define instead
\[
 \gamma(F)=\mathop{\operatorname{colim}}_{\varnothing\ne U\subseteq T}F(U),
\]
where transition maps are restrictions as $U$ shrinks. Nonempty opens form a filtering system under intersection. If a generic point exists, this is its stalk; otherwise it is the stalk at the corresponding point of the soberification. The following direct argument avoids requiring that point.

For every nonempty open $U$, a section of $F(U)$ has a well-defined generic germ. Because $U$ is connected, the constant sheaf with value $\gamma(F)$ has sections $\gamma(F)$ on $U$. We obtain a map $F\to\underline{\gamma(F)}$ and a decomposition
\[
 F\cong\coprod_{g\in\gamma(F)}F_g
\]
into its fibres. Every $F_g$ is noninitial because $g$ has a representative section. If $A,B$ are two noninitial subobjects of $F_g$, choose sections $a\in A(U)$ and $b\in B(V)$ on nonempty opens. They have the same generic germ, so on some nonempty $W\subseteq U\cap V$ their restrictions agree. This supplies a section of $A\cap B$, proving that $F_g$ is irreducible.

Consequently these fibres are exactly the connected components of $F$, and $F$ is $0$-connected if and only if $\gamma(F)$ is a singleton. This proves both the assertion about the generic fibre and the missing implication in footnote~7: the sheaf topos of an irreducible space is locally irreducible. More generally, if $T$ is a disjoint union of open irreducible spaces, the conclusion follows component by component. Conversely, a locally irreducible sheaf topos is locally connected; its terminal object decomposes into connected subterminal objects, which are disjoint open subsets of $T$, and each is irreducible. This proves the topological characterization stated here.

For a general presheaf topos, no generic point of the underlying preorder can replace this argument. Parallel arrows produce different objects in slices, even when they have the same source and target. The next example shows exactly where the analogy with topological spaces ceases to suffice.
\end{dcomment}

\begin{dcomment}{exposition}{p15:J'aimerais une condition}{Completing Proposition 2.1}
Proposition~2.1 isolates the part of filtering which can be expressed without equalizing parallel arrows. We give its proof because it explains both the corrected statement and the mechanism which will later be strengthened.

Recall that PS~1 says that every cospan $a\to x\leftarrow b$ admits a commuting square with some upper-left object $c$; no universal property is requested. Suppose PS~1 holds. In any category with this property, the relation ``$a$ and $b$ have a common predecessor'' is reflexive and symmetric. It is transitive: given $a\leftarrow r\to b$ and $b\leftarrow s\to d$, apply PS~1 to $r\to b\leftarrow s$ and compose the resulting arrows. Every arrow connects objects related in this sense. Hence any two objects in the same connected component have a common predecessor.

For a presheaf $F$, its category of elements $E=\int_XF\to X$ is a discrete fibration, the meaning of ``\'etale'' here. Every arrow into the image of an object has a unique lift into that object. Thus a PS~1 square in $X$ lifts to a PS~1 square in $E$: the two lifts into the common target agree by uniqueness. Therefore $E$ also satisfies PS~1. Components of $E$ correspond to coproduct components of $F$. If $F$ is $0$-connected, $E$ is nonempty and connected, so the previous paragraph gives common predecessors for every pair of its objects. Subobjects of $F$ correspond to sieves in $E$; any two nonempty such sieves meet. Thus $F$ is irreducible, and $\widehat X$ is locally irreducible.

Conversely, suppose every $0$-connected presheaf is irreducible. The representable $h_x$ is $0$-connected, since its category of elements $X/x$ has a terminal object. Given arrows $a\to x$ and $b\to x$, their generated subobjects of $h_x$ are noninitial. Their intersection is noninitial, so it contains some element represented by arrows $c\to a$ and $c\to b$ with equal composite to $x$. This is PS~1. The descriptions using $X/x$, its associated space, and open subcategories are equivalent by the preceding sieve argument. This proves all the equivalences of Proposition~2.1.

The counterexample is the category with two objects and distinct arrows $\alpha,\beta:0\rightrightarrows1$. Its objects have the common predecessor $0$, so its underlying preorder is the irreducible two-point chain. But in $X/1$, the objects $(0,\alpha)$ and $(0,\beta)$ have no common predecessor: the only endomorphism of $0$ is its identity. Their two singleton lower sets are disjoint nonempty subobjects of $h_1$. Thus $\widehat X$ is not locally irreducible. The mistaken converse on page~16 is precisely the assertion that the preorder remembers this equality of composites; the footnotes already retract it.
\end{dcomment}

\begin{dcomment}{examples}{p16:mais la réciproque est fausse.}{The equalization condition survives beyond every slice}
The definitions now separate PS~1 from the additional condition PS~2. In the latter, ``\'egalisateur'' means an \emph{equalizing arrow}: for $u,v:a\rightrightarrows b$ there is $w:c\to a$ with $uw=vw$. It does not mean that a categorical equalizer exists. This is the usual cofiltered condition, dual to SGA~4's PS~2 \citep[Expos\'e~I, \S2.7]{SGA4I1972}.

To see why all slices can be cofiltered without the original category being locally cofiltered, take the one-object category $BG$ for a nontrivial group $G$. It satisfies PS~1: a pair $f,g$ can be completed by $1$ and $g^{-1}f$. For its unique object $*$, the slice $BG/*$ has objects $g\in G$ and exactly one arrow from $g$ to $h$, namely $h^{-1}g$. It is equivalent to the terminal category and is cofiltered. Yet distinct $u,v\in G$ cannot be equalized on the right, since $uw=vw$ and invertibility of $w$ imply $u=v$. Hence $BG$ fails PS~2.

This example also makes the topos distinction visible. Every connected $G$-set is transitive and therefore irreducible, so the presheaf topos is locally irreducible. Products of connected objects need not remain connected: for the regular right $G$-set, the diagonal right action on $G\times G$ has one orbit for each value of $hg^{-1}$. The product condition introduced next detects the missing equalization information. No defect in Proposition~2.1 is involved; its purpose was exactly to isolate PS~1.
\end{dcomment}


% Manuscript pages 17--35. Commentary prepared for consolidation.

\begin{dcomment}{development}{p16:avec $(\alpha\alpha' = \beta\beta')}{The finite example and its homotopy type}
The example between Definition 2.2 and Proposition 2.3 makes the loss of information in $\operatorname{Ord}(X)$ concrete. We read the displayed graph and its two relations as a presentation of $X$; in particular, no additional relation is imposed on composites with $\gamma$. Put
\[
 r=\alpha\alpha'=\beta\beta',\qquad
 s=\alpha\beta'=\beta\alpha',\qquad r\ne s.
\]
Then $\operatorname{Ord}(X)$ is the chain $0<1<2<3$. Every pair of objects has both a common predecessor and a common successor. Nevertheless, $0$ is not initial: the arrows $\alpha'\gamma,\beta'\gamma:0\rightrightarrows2$ are distinct. Likewise, $3$ is not terminal, since $r,s:1\rightrightarrows3$ are distinct.

The same calculation locates the exact failure of PS 1. A commuting square for the cospan $1\xrightarrow{\alpha'}2\xleftarrow{\beta'}1$ would require two arrows from an object into $1$. Their common source can only be $1$ or $0$, and both arrows would be respectively $1_1$ or $\gamma$. Neither choice makes the square commute. Thus the left PS 1 fails, and the pair $\alpha',\beta'$ also has no equalizing arrow on the left. Dually, a commuting cocone over $3\xleftarrow{\alpha}2\xrightarrow{\beta}3$ can only use identity arrows out of $3$, so the right PS 1 fails; the pair $\alpha,\beta$ cannot be equalized on the right either. The category is consequently neither locally cofiltered nor locally filtered, although its associated poset satisfies both conditions. This is the point of the paragraph immediately following the example.

The crossed composition relations also admit a useful homotopical interpretation which the manuscript does not develop.

\textbf{The homotopy type of the finite example.}\ours
\textit{For the category presented above, $|N X|\simeq\mathbb{R}P^2$. In particular, its contractible reachability poset does not determine even the fundamental group of its nerve.}

\textit{Proof.} Let $C$ be the full subcategory on $1,2,3$. The functor $q:X\to C$ sends $0$ to $1$, sends $\gamma$ to $1_1$, and fixes the remaining generators. It respects the defining relations. For the inclusion $j:C\hookrightarrow X$, one has $qj=1_C$ and a natural transformation $1_X\to jq$ whose component at $0$ is $\gamma$ and whose other components are identities. The resulting homotopy of nerves fixes $N C$, so $|N X|\simeq|N C|$.

Identify each set $C(i,j)$, for $1\le i<j\le3$, with $\mathbf Z/2$: label $\alpha',\alpha,r$ by $0$ and $\beta',\beta,s$ by $1$. The two displayed relations say exactly that composition adds these labels. Let $P$ be the poset with objects $(i,\varepsilon)$, where $i\in\{1,2,3\}$ and $\varepsilon\in\mathbf Z/2$, with $(i,\varepsilon)<(j,\delta)$ precisely when $i<j$. Its involution changes every $\varepsilon$ to $\varepsilon+1$. There is an invariant functor $P\to C$ sending $(i,\varepsilon)$ to $i$ and an arrow to the label $\delta-\varepsilon$. A chain of arrows in $C$ lifts uniquely after choosing the lift of its first vertex, and the two choices are exchanged by the involution. This proves, degree by degree and compatibly with faces and degeneracies,
\[
 N P/(\mathbf Z/2)\simeq N C.
\]
The order complex of $P$ is the join of three two-point sets. Sending their vertices to $\{\pm e_1,\pm e_2,\pm e_3\}$ identifies it with the boundary of the octahedron. Radial projection identifies its realization with $S^2$, and the involution with the antipodal map. Realization commutes with the quotient, giving
\[
 |N C|\simeq |N P|/(\mathbf Z/2)\cong S^2/(x\sim-x)=\mathbb{R}P^2.
\]
Since the sphere is simply connected, this double covering also gives $\pi _1(|N X|)\cong\mathbf Z/2$.\hfill$\square$

The join and antipodal descriptions are classical \citep[p.~9 and Example 1.43]{Hatcher2002}. Our addition is their explicit identification with this particular presentation. The bounded literature searches did not locate that identification; the result is marked as our calculation, without a claim of priority.
\end{dcomment}

\begin{dcomment}{exposition}{p17:Reste à exprimer la condition filtrante elle-même}{Products of connected objects}
Grothendieck now asks for a condition on the topos which remembers the equalization of parallel arrows. Irreducibility alone cannot do this: it records whether two subobjects meet, whereas equalization concerns two maps with the same source. Proposition 2.3 first isolates the intermediate condition that the component functor preserve products. Throughout the following topos discussion, ``connected object'' means noninitial and connected, that is, $0$-connected; the topos is locally connected. Thus every object is the coproduct of its connected components, and $\pi _0$ is left adjoint to the constant-object functor.

Write $\mathcal E$ for the category of objects of the topos and let $S$ be the full generating subcategory of connected objects used in the text. The precise statement of 2.3(b) is
\[
 \pi _0(V\times U)\longrightarrow\pi _0(V)
 \quad\hbox{is bijective whenever $U$ is connected.}
\]
The edition prints $\pi _0(U)$ on the right, although the displayed projection has target $V$. This is a slip of the edition, retained in the source. Condition (a) implies (c) by decomposing $U$ and $V$ into connected components: products distribute over coproducts in a topos, and each resulting product is connected. Conversely, (c) applied to connected objects gives (a); the same decomposition proves the equivalence with (b).

The reduction to generators deserves the proof which the footnote requests. For $F\in\mathcal E$, let $S/F$ have objects $s\to F$, with $s\in S$. The canonical presentation
\[
 F\simeq\mathop{\rm colim}_{(s\to F)\in S/F}s
\]
holds in $\mathcal E$; it is the density presentation for the canonical site of generators. Applying $\pi _0$ gives $\pi _0(F)\simeq\pi _0(S/F)$, since every $s$ is connected. If $s\times t$ is connected for every $s,t\in S$, separate preservation of colimits by products gives
\[
 \pi _0(F\times G)
 \simeq\mathop{\rm colim}_{S/F\times S/G}\{*\}
 \simeq\pi _0(S/F)\times\pi _0(S/G).
\]
This proves (d)$\Rightarrow$(c). Equivalently, one may use the clopen-partition argument in our block on Lemma 2.9 below; this avoids any choice of a presentation.

Applying (c) to the terminal object shows that the diagonal of $\pi _0(1)$ is a bijection onto its square. Consequently $\pi _0(1)$ has at most one element. The nondegenerate case has exactly one, which explains the separate nonemptiness clause in Definition 2.4. Nothing in this argument annihilates higher homotopy groups. The inference to $W_\infty$-asphericity immediately following 2.3 is explicitly withdrawn in the manuscript's annotation; the stronger hereditary hypothesis introduced in 2.6 is what will support the cohomological argument.
\end{dcomment}

\begin{dcomment}{later}{p18:Notons que pour un objet}{The ordinary sifted condition and its variance}
The product condition has acquired a useful name, provided that we keep track of the opposite category. For $\widehat X=\operatorname{Fun}(X^{\mathrm{op}},\mathbf {Set})$, the left adjoint to the constant-presheaf functor is
\[
 \pi _0(F)=\mathop{\rm colim}_{X^{\mathrm{op}}}F
          =\pi _0(X/F).
\]
Thus $\widehat X$ is totally $0$-connected precisely when $X^{\mathrm{op}}$ is \emph{sifted}: its colimits in sets commute with finite products. In the ordinary categorical characterization, the indexing category is nonempty and its diagonal is final. For the conventions of these pages, this says that $X\to X\times X$ is co-cofinal. The comma category at $(x,y)$ consists of common cones $z\to x$, $z\to y$, and is exactly $X/(h_x\times h_y)$.

This characterization also explains why binary products suffice. Repeated use of finality of the diagonal gives finality of every finite nonempty diagonal; nonemptiness handles the nullary product. The terminology and its relation to algebraic theories are developed by \citet{AdamekRosickyVitale2010}. The distinction from filteredness is structural: filtered colimits preserve finite limits, while sifted colimits are required to preserve finite products. Equalizers are precisely the operation absent from the latter requirement.

There is a corresponding distinction in higher category theory, but one must not transfer a statement by changing its vocabulary alone. In the higher version the relevant comma objects must be weakly contractible, rather than merely nonempty and connected. The characterization of filtered colimits by preservation of finite limits in spaces is \citet[Proposition 5.3.3.3]{Lurie2009HTT}; sifted colimits and geometric realizations are treated in his \S5.5.8. Our present identification concerns \emph{ordinary} set-valued colimits. The higher homotopy obstruction which survives it is exhibited later in the commentary on 2.19--2.20.
\end{dcomment}

\begin{dcomment}{examples}{p19:les catégories induites}{The simplex category separates the two conditions}
Grothendieck's example shows exactly where a global product condition loses the information needed locally. We spell it out because the word ``visibly'' hides two different computations. Let $\Delta$ have the nonempty finite ordinals $[n]=\{0,\ldots,n\}$ and monotone maps. For any small category $X$ with binary products, the product of two representables is representable; its category of elements has a terminal object. Such a nonempty $X$ is therefore totally $0$-connected. The separate assertion about $\Delta$ does not say that $\Delta$ has binary products: the ordered square $[1]\times[1]$ is not an ordinal.

For $\Delta$, the product $h_{[m]}\times h_{[n]}$ is the simplicial set $\Delta^m\times\Delta^n$. Its vertices are the pairs $(i,j)$. An edge joins $(i,j)$ to $(i+1,j)$ whenever the latter exists, and similarly joins $(i,j)$ to $(i,j+1)$. Hence its vertex graph is connected. Every simplex is joined in its category of elements to each of its vertices by a face map. The category of elements is consequently nonempty and connected, which verifies the generator test of 2.3 for all $m,n$.

Local irreducibility fails already over $[1]$. In $\Delta/[1]$, consider the two objects $[0]\xrightarrow{0}[1]$ and $[0]\xrightarrow{1}[1]$. A common predecessor would be a nonempty ordinal $[r]$ mapping to both copies of $[0]$, with equal composites to $[1]$. Those composites are respectively constant at $0$ and $1$, so equality is impossible. The two principal sieves have empty intersection. In particular, $\Delta/[1]$ is not irreducible, and $\Delta$ is not cofiltered. Equivalently, the parallel vertex maps $[0]\rightrightarrows[1]$ admit no equalizing map into $[0]$.

This example gives an explicit positive answer to the request in footnote 31: total $0$-connectedness does not imply local irreducibility. Its opposite also supplies an ordinary sifted category which is not filtered. Neither conclusion requires an argument about derived limits.
\end{dcomment}

\begin{dcomment}{exposition}{p20:Montrons quand même les deux implications}{What the hereditary condition adds}
Theorem 2.6 makes the product condition stable under passage to a connected object as base. Its surviving equivalence is
\[
 (a)\ \Longleftrightarrow\ (a')\ \Longleftrightarrow\ (b)
 \ \Longleftrightarrow\ (d),
 \qquad (c)\Longleftrightarrow(c'),
\]
with an implication from the first group to the second. Here the three objects $U,V,W$ in (a') are all connected. Connectedness of $W$ is essential: even in sets, two singleton objects mapping to different elements of a two-element set have empty pullback.

The equivalence of (a) and (a') is the same product test inside $\mathcal E/W$. Condition (a) implies local irreducibility by 2.5 and gives the assertion about components in (b). To see the decisive step (b)$\Rightarrow$(d), first use 2.8 for PS 1. Given $u,v:a\rightrightarrows b$ in $S$, both objects lie over the same connected component of $1$. By (b), $b\times b$ is connected. The equalizer is the pullback
\[
 E=a\times_{b\times b}b,
 \qquad a\longrightarrow b\times b=(u,v),\quad
 b\longrightarrow b\times b=\operatorname{diag}.
\]
Local irreducibility, applied over the connected object $b\times b$, makes $E$ noninitial. A generator mapping to $E$ supplies an equalizing arrow, which is PS 2. The text's ``c) est vérifié'' at manuscript p.~24 means (b); its reference to $\mathcal X\times\mathcal X$ should be read as an application of 2.8 in $\mathcal X$ itself.

There is a concise exactness interpretation after restricting to any connected component of the topos. Under (a'), decomposition into connected components gives, for arbitrary $U,V,W$ in that component,
\[
 \pi _0(U\times_WV)
 \simeq\pi _0(U)\times_{\pi _0(W)}\pi _0(V).
\]
Indeed, the summands of the pullback are precisely the products $U_i\times_{W_k}V_j$ for which both factors map to the same component $W_k$; each is connected by (a'). Since $\pi _0(1)=\{*\}$, $\pi _0$ preserves finite limits on this connected topos. We shall use this stronger fact in the cohomological argument. For a disconnected topos one applies it component by component; its global $\pi _0$ does not preserve the terminal object.
\end{dcomment}

\begin{dcomment}{exposition}{p23:On a prouvé 2.8, donc aussi}{How irreducible generators control an arbitrary slice}
The proof of 2.8 explains why irreducibility, although weaker than filteredness, is genuinely local in the topos. The missing hypothesis in its printed condition (b) is that $W$ be connected. The assertion used in the proof is (b'): noninitial $U$ and $V$ over a connected $W$ have noninitial pullback.

Here is the passage from the generating objects to such a $W$. Suppose every $s\in S$ is irreducible, and consider the images $W_s\subseteq W$ of all maps $s\to W$. These images cover $W$ and remain irreducible: inverse images of two nonzero subobjects under an epimorphism are nonzero, and their nonzero intersection has nonzero image. On this covering family, declare two images related when they meet. The relation is reflexive and symmetric. It is transitive because, if $A$ and $C$ both meet an irreducible $B$, the two nonzero subobjects $A\cap B$ and $B\cap C$ of $B$ meet. Thus the equivalence classes give pairwise disjoint unions of subobjects covering $W$.

These unions are a coproduct decomposition in the topos. Since $W$ is connected, there is only one class. Every pair of generating images therefore meets. Choose a generator mapping into each of the noninitial objects $U$ and $V$. Their images in $W$ meet, and epimorphisms remain epimorphisms after pullback; consequently $U\times_WV$ is noninitial. This proves the step (c)$\Rightarrow$(b') without assuming that every object of a site is representable as a presheaf on that site.

Finally, the condition that $s\times_x t$ be noninitial is equivalent to the existence of a generator $r$ mapping to that pullback. Such a map is exactly a commutative square completing $s\to x\leftarrow t$, namely PS 1. This proves the equivalence with (d) and (e). The implication (b')$\Rightarrow$(a) follows by taking two nonzero subobjects of a connected object. The remaining implications are restrictions of this same test. Grothendieck's argument is therefore sound after restoring the condition on the base; the diagram of implications is not concealing an additional homotopical assumption.
\end{dcomment}

\begin{dcomment}{exposition}{p25:Je n'explicite pas la démonstration}{Connectedness in a category of generators}
Lemma 2.9 is what permits the categorical and topos-theoretic arguments to meet. Its content is stronger than the statement that a generating family covers the terminal object: a partition of the generators must descend to an actual coproduct decomposition of that object.

Let $S\subseteq\mathcal E$ be full, generating, and formed of connected objects. A decomposition $1=A\amalg B$ assigns every $s\in S$ entirely to one summand, since the two pullbacks decompose $s$. Every arrow preserves this assignment. It thus gives a partition $S=S_A\amalg S_B$ into full subcategories with no arrows between them. Conversely, suppose such a partition is given. Let $A$ and $B$ be the unions in $1$ of the images of the objects in the respective parts. They cover $1$, because $S$ generates. If $A\cap B$ were noninitial, some generator would map into it. Pulling back the epimorphic families defining $A$ and $B$, and choosing a generator in a noninitial member, would produce an object $t\in S$ with maps to an object of $S_A$ and one of $S_B$. This contradicts the partition. Hence $A\cap B=0$, and $A\amalg B\simeq1$.

The two constructions are inverse: a connected generator whose image lies in $A$ cannot belong to the other part, and the images of all generators assigned to $A$ jointly cover $A$. The same argument works with any set of parts and proves
\[
 \pi _0(S)\simeq\pi _0(1_{\mathcal E}).
\]
It also applies to $S/F$, which is a full generating subcategory of $\mathcal E/F$. An object $s\to F$ is connected in the slice exactly when its domain $s$ is connected, since coproducts in a slice are computed on domains. We obtain the repeatedly used identity
\[
 \pi _0(S/F)\simeq\pi _0(F).
\]
This supplies the proof omitted at manuscript p.~25 and justifies the earlier use of the lemma in 2.8. It also explains why the category of elements, rather than only the set of objects of a site, controls connectedness.
\end{dcomment}

\begin{dcomment}{exposition}{p27:mais pas l'implication inverse}{Where both equalizations are used}
The completion of 2.6 is a useful instance of the finite-cone argument underlying filteredness. It is not enough to find a common predecessor of two objects: one must also arrange that the resulting maps agree in both factors. This is why the proof of Lemma 2.10 invokes PS 2 twice.

Suppose $S$ is cofiltered. An object of $S/(h_a\times h_b)$ is a triple $(x,r:x\to a,s:x\to b)$. It exists by the common-predecessor axiom. Given two such triples with domains $x$ and $y$, choose $z\to x$ and $z\to y$. Equalize the two resulting maps $z\rightrightarrows a$ by a map $z'\to z$, then equalize the two maps $z'\rightrightarrows b$ by $z''\to z'$. The second operation preserves the first equality. The object with domain $z''$ is a common predecessor in $S/(h_a\times h_b)$. This proves its connectedness, as required in 2.10. In fact it is cofiltered: parallel arrows in it can be equalized after projection to $S$, and the resulting arrow still belongs to the category of elements.

Now let $S$ merely be locally cofiltered, and let $W$ be connected in the topos. The projection $S/W\to S$ is a discrete fibration. PS 1 lifts along a discrete fibration: once a common predecessor is chosen downstairs, lifting one leg determines the object upstairs, and equality of the two composites gives the other leg. PS 2 lifts for the same reason, using uniqueness of lifts. Thus $S/W$ is locally cofiltered. Lemma 2.9 says it is connected, so it is cofiltered. Applying the preceding argument in this generating category and then the generator test of 2.3 shows that $\mathcal E/W$ is totally $0$-connected. This is (d)$\Rightarrow$(a).

The proof establishes the announced equivalence of (a), (a'), (b), and (d). It does not establish (c')$\Rightarrow$PS 2, and the annotation correctly removes that assertion. Restricting the base of a pullback to a representable object has lost information: the object $b\times b$ required for the diagonal argument need not be among the chosen generators. The following group example makes this loss explicit.
\end{dcomment}

\begin{dcomment}{examples}{p28:Mais la réciproque est-elle vraie}{A groupoid invisible to its representable slices}
The marginal answer to Perplexity 2.11 is decisive, and it also explains the weaker condition that remains. Let $G$ be a nontrivial group and let $BG$ be its one-object category. The objects of $BG/*$ are the elements $g\in G$, viewed as arrows to the unique object. For any $g,h$, there is exactly one arrow from $g$ to $h$, namely the element $k$ satisfying $hk=g$. Hence $BG/*$ is a nonempty indiscrete groupoid, equivalent to the terminal category, and is cofiltered.

The category $BG$ itself is not cofiltered. If $g\ne h$ are parallel arrows, an equalization $gk=hk$ would imply $g=h$ because $k$ is invertible. Its presheaf topos is the category of right $G$-sets. The representable is the regular $G$-set $G$; it is connected. Its square with diagonal action has
\[
 \pi _0(G\times G)\simeq G,
 \qquad [(a,b)]\longmapsto ba^{-1},
\]
after choosing the corresponding convention for the right action. Thus the square is not connected. This verifies directly that the topos is locally irreducible but is not totally $0$-connected, although the slice over the representable is equivalent to sets. The image of the regular object onto the terminal $G$-set is exactly the quotient missed by the generator-only test.

There is, however, no further example in which \emph{every} $X/x$ is totally $0$-connected but some $X/x$ is not cofiltered. Indeed, total $0$-connectedness of all $X/x$ makes all these categories irreducible, so $X$ is locally irreducible. For $y\to x$, the slice $(X/x)/(y\to x)$ is isomorphic to $X/y$, and is therefore irreducible. Hence $X/x$ is both totally $0$-connected and locally irreducible. Corollary 2.7, whose proof has just been completed, makes it cofiltered.

In elementary terms, the surviving weaker condition is PS 1 together with
\[
 a\mathrel{\substack{\xrightarrow{u}\\[-.5ex]\xrightarrow[v]{}}}b
 \xrightarrow{h}x,\quad hu=hv
 \quad\Longrightarrow\quad
 \exists\,t:c\to a\ \text{with }ut=vt.
\]
This is PS$'$ 2 in footnote 35. It equalizes parallel arrows only when they already become equal after a common outgoing arrow. In a groupoid that antecedent forces equality, so the condition is automatic. PS 2 has no such antecedent. The two questions in the paragraph are therefore settled differently: the first has the groupoid counterexample already supplied by Grothendieck; the second has a negative existence answer from his preceding theorem.
\end{dcomment}

\begin{dcomment}{exposition}{p29:C'est la forme `topologique' la plus faible}{Why the square of one generator detects equalizers}
Corollary 2.12 is not merely another presentation of a long list of equivalent conditions. It isolates the diagonal as the part of the product test which forces equalization. We can read its criterion entirely in the topos, without treating the chosen site as a space.

Suppose that $s\times s$ is irreducible for every $s\in S$. The diagonal $s\hookrightarrow s\times s$ is a noninitial subobject. Every noninitial subobject of an irreducible object is irreducible, so each $s$ is irreducible and 2.8 gives PS 1. For two maps $u,v:a\rightrightarrows b$ in $S$, the image of $(u,v):a\to b\times b$ is noninitial. It meets the diagonal, because $b\times b$ is irreducible. Pullback stability of epimorphisms then says that
\[
 \operatorname{Eq}(u,v)=a\times_{b\times b}b
\]
is noninitial. A generator mapping into this equalizer supplies PS 2. Conversely, PS 1 and PS 2 give local cofilteredness; the category of generators over $b\times b$ is cofiltered by the proof of 2.10, so this product is irreducible.

This also proves the comment that the separate hypothesis on $\mathcal E/s$ is redundant. The diagonal is a nonzero open subobject in the language of slice topoi, and irreducibility passes to such an open. Saying instead that $s$ is irreducible and its diagonal is dense gives the same test: density means that the diagonal meets the image of every noninitial object mapping to $s\times s$. The pair $(u,v)$ then gives the equalizer test just used. The relevant ``open immersion'' is the morphism of topoi associated with a monomorphism in $\mathcal E$; it is not an assertion that an arbitrary topological diagonal is open.

Definition 2.13 packages these equivalent conditions as local total $0$-connectedness. The generator category is allowed to vary: once the condition is expressed for all connected slices of the topos, the proof of 2.6 applies to every full generating subcategory of connected objects.
\end{dcomment}

\begin{dcomment}{examples}{p30:la cohomologie d'un espace irréductible est triviale}{Irreducible spaces still have nonconstant cohomology}
The reference to Serre separates two coefficient questions which must remain distinct. On an irreducible nonempty topological space $T$, every nonempty open is connected. For an abelian group $M$, the constant sheaf $\underline M$ consequently has sections $M$ on every nonempty open, with identity restriction maps. Restrictions to the empty open are also surjective. Thus $\underline M$ is flabby and
\[
 H^q(T,\underline M)=0\quad(q>0),\qquad
 H^0(T,\underline M)=M.
\]
The same argument on each irreducible component applies when the space is a disjoint union of irreducible open components, which is the local condition used in 2.14.

For a concrete contrast, take the Zariski space $T=\mathbf P^1_k$, which is irreducible, and the line bundle $\mathcal O(-2)$. On the standard affine charts with coordinate $t$ and $t^{-1}$, its degree-one \v{C}ech quotient is
\[
 \frac{k[t,t^{-1}]}{k[t]+t^{-2}k[t^{-1}]}
   =k\,[t^{-1}].
\]
The first summand removes all nonnegative powers and the second all powers at most $-2$; exactly the coefficient of $t^{-1}$ remains. The affine charts and their intersection are acyclic for this coherent sheaf, so the quotient is $H^1(\mathbf P^1_k,\mathcal O(-2))\simeq k$. The affine-cover comparison used here is \citet[no.~47, Theorem 4]{Serre1955FAC}; the displayed computation gives the particular example in full.

For 2.14 itself, choose the connected nonempty opens as generators. If the space is locally irreducible in Grothendieck's sense, each such open is irreducible. Two nonempty opens inside it meet, giving PS 1; the generator category is a poset, so parallel arrows are already equal and PS 2 imposes nothing further. The converse follows from local total $0$-connectedness implying local irreducibility. This is why the two notions coincide for spaces while the group example separates them for general topoi.
\end{dcomment}

\begin{dcomment}{exposition}{p31:Supposons que ce soit prouvé pour les}{The hereditary hypothesis in the cohomological proof}
Grothendieck's induction proves the intended assertion after keeping its hypothesis fixed: the topos must be connected and \emph{locally} totally $0$-connected. The word ``localement'' disappears in the sentence introducing the induction at manuscript p.~31, but it is needed when the argument passes to components of powers of a covering object. Those components define connected slices of the original topos, and hence remain in the same class by 2.6. Mere total $0$-connectedness does not provide this inheritance.

First consider a connected locally constant object $L$. Its diagonal in $L\times L$ is complemented, as one checks after a cover trivializing $L$. The product is connected, so this noninitial summand is the whole product. Hence $L\to1$ is monic. Its image is also complemented, by the same local triviality; connectedness of the base and nonemptiness of $L$ give $L\simeq1$. Every covering object is therefore a coproduct of copies of $1$. In the source paragraph, $U$ must still denote a connected component of a \emph{covering}, not an arbitrary connected sheaf.

For the higher cohomology, assume the vanishing in degrees $0<q<n$ for every connected topos in this hereditary class. Given $\xi\in H^n(\mathcal E,\underline M)$, choose an epimorphism $U\to1$ killing $\xi$. Such a cover exists because a positive-degree cohomology class represented in an injective resolution becomes a boundary locally. The augmented \v{C}ech nerve gives the descent spectral sequence
\[
 E_1^{pq}=H^q(\mathcal E/U^{p+1},\underline M)
 \ \Longrightarrow\ H^{p+q}(\mathcal E,\underline M).
\]
Each $U^{p+1}$ is a coproduct of connected objects whose slices remain in the class; thus the rows $0<q<n$ vanish. The edge sequence gives
\[
 0\longrightarrow H^n(C^\bullet)
 \longrightarrow H^n(\mathcal E,\underline M)
 \longrightarrow \ker\bigl(H^n(U,\underline M)
       \rightrightarrows H^n(U^2,\underline M)\bigr).
\]
The last image of $\xi$ is zero by choice of $U$. If $I=\pi _0(U)$, product preservation identifies
\[
 C^p=\operatorname{Map}(I^{p+1},M).
\]
Choose $i_*\in I$, which exists since $U\to1$ is epic. Inserting $i_*$ as first coordinate defines
$(sc)(i_0,\ldots,i_{p-1})=c(i_*,i_0,\ldots,i_{p-1})$.
For the alternating deletion differential, direct cancellation gives $ds+sd=1$ in positive degrees, with the usual augmentation in degree zero. Thus $H^n(C^\bullet)=0$ and $\xi=0$. The ``trivial'' complex in the manuscript is this contractible augmented complex; its individual groups and differentials are not zero.
\end{dcomment}

\begin{dcomment}{later}{p32:en degrés $>0$ est nulle}{From connected components to the hypercover homotopy type}
The stronger homotopical conclusion of 2.6 can also be proved directly, without deducing weak contractibility from an unspecified version of a cohomological Whitehead theorem. On a connected locally totally $0$-connected topos, the component functor preserves finite limits, as shown above, and preserves epimorphisms: epimorphisms in a topos are regular, and the left adjoint $\pi _0$ preserves their coequalizer presentations. Let $U_\bullet\to1$ be a hypercover by objects of this topos. Its matching objects are finite limits. Therefore
\[
 \pi _0\bigl(M_n(U_\bullet)\bigr)
       \simeq M_n\bigl(\pi _0(U_\bullet)\bigr),
\]
and the hypercover epimorphisms become surjections of sets. The simplicial set $\pi _0(U_\bullet)$ consequently has a filler for every boundary $\partial\Delta^n\to\pi _0(U_\bullet)$, including the nonemptiness condition in degree zero. Its map to a point is a trivial Kan fibration, and it is weakly contractible.

In the hypercover construction of the homotopy type, whose classical setting is \citet[\S\S8--9]{ArtinMazur1969}, every stage is therefore weakly contractible. This proves the $W_\infty$-asphericity meant here, as well as the corresponding assertion for each connected object by passage to its slice. For a presheaf topos it says concretely that every connected category of elements is weakly contractible.

There is also a short algebraic verification of the constant-coefficient calculation. Since $\pi _0$ preserves finite limits, it induces an exact functor on abelian group objects, left adjoint to the constant-abelian-sheaf functor. The latter therefore preserves injectives. Applying it to an injective resolution of $M$, and then taking global sections, recovers the same resolution of $M$, because the topos is connected. Thus $R\Gamma(\mathcal E,\underline M)\simeq M[0]$. These proofs make explicit why the hereditary condition repairs the false inference after 2.3. They impose no claim that the topos itself is equivalent to sets; the preceding projective-line example shows that arbitrary coefficients retain substantial information.
\end{dcomment}

\begin{dcomment}{exposition}{p32:C'est déjà dans SGA 4}{The terminal presheaf as an ind-object}
The passage to ind-objects replaces the cone conditions by a representability statement. The general form, recorded in 2.18, is the familiar criterion from \citet[Expos\'e I, Theorem 8.3.3]{SGA4I1972}: a presheaf $F$ on a small category $X$ is a filtered colimit of representables if and only if its category of elements $X/F$ is filtered. We give the argument, since it also supplies all the details of 2.15.

For every presheaf there is the density presentation
\[
 F\simeq\mathop{\rm colim}_{(x,s)\in X/F}h_x.
\]
Objectwise, the element represented by a pair $a\to x$, $s\in F(x)$ maps to its pullback in $F(a)$; it has the representative $(a,t)$ with identity map, and the relations identify every other representative with this one. If $X/F$ is filtered, this is the required ind-presentation.

Conversely, write $F=\operatorname{colim}_{i\in I}h_{x_i}$ with $I$ filtered. The category of elements is nonempty, since the image of $1_{x_i}$ supplies an object. Two elements $s\in F(a)$ and $t\in F(b)$ have representatives at stages $i,j$; passage to a common stage $k$ supplies arrows from $(a,s)$ and $(b,t)$ to $(x_k,\overline{1}_{x_k})$. Given parallel arrows $u,v:(a,s)\rightrightarrows(b,t)$, represent $t$ by $b\to x_i$. The two composites $a\rightrightarrows x_i$ have the same image in $F(a)$. Equality in a filtered colimit of sets is witnessed at a common later stage. The resulting map $(b,t)\to(x_j,\overline{1}_{x_j})$ equalizes $u,v$. Thus $X/F$ is filtered.

For $F=1$, its category of elements is $X$ itself, giving 2.15. The formula $1=\operatorname{colim}_{x\in X}h_x$ is valid even when $X$ is not filtered: at $a$ its value is $\pi _0(a\backslash X)=\{*\}$ because $(a,1_a)$ is initial. The condition lies in the possibility of a \emph{filtered} presentation. The reindexing by a directed poset used in the manuscript is Deligne's result, \citet[Expos\'e I, Proposition 8.1.6]{SGA4I1972}; the proof above does not need this extra reindexing.
\end{dcomment}

\begin{dcomment}{later}{p34:Plus généralement et agréablement}{Filtered categories and the minimal localizer}
Corollary 2.16 needs only the standard contractibility of a filtered category, together with the comparison between its weak equivalences and the chosen localizer. Here is the finite-data argument behind the first assertion. In a filtered category, any finite collection of objects has maps into one object. Given finitely many arrows between those objects, successive applications of the parallel-arrow axiom modify the common target so that all the required triangles commute. Thus every finite diagram of objects and arrows occurring in a finite simplicial subset of the nerve admits a compatible cone.

Every map from a finite sphere into the realization of the nerve can, after simplicial approximation, be represented in a finite such subset. The cone makes that map nullhomotopic. The nerve is nonempty and connected, so all its homotopy groups vanish. Reversing the category does not change the homotopy type of its nerve, and gives the cofiltered assertion as well.

The manuscript formulates the conclusion for a fundamental localizer satisfying its axiom Loc(8). A later theorem permits a sharper modern statement: the nerve weak equivalences $W_\infty$ form the smallest fundamental localizer. This is \citet[Theorem 2.2.11, p.~136]{Cisinski2004Minimal}. With the standard modern definition of fundamental localizer, every filtered or cofiltered category is consequently $W$-aspherical without an additional filtered-limit axiom. This is a later strengthening of the hypothesis in the manuscript, rather than a reason to alter its historical formulation. The finite-cone argument also underlies the higher categorical characterization cited earlier: filtered colimits of spaces commute with finite limits.
\end{dcomment}

\begin{dcomment}{exposition}{p35:Donc finalement, on trouve des exemples non triviaux}{Left Kan extension carries the ind-presentation}
The reason behind 2.17 is that an ordinary left Kan extension sends a representable to a representable and preserves colimits. If $f:X\to Y$ and $X$ is filtered, then
\[
 f_!1_X\simeq\mathop{\rm colim}_{x\in X}f_!h_x
          \simeq\mathop{\rm colim}_{x\in X}h_{f(x)}.
\]
The result is an ind-representable presheaf on $Y$, so $Y/(f_!1_X)$ is filtered by the criterion just proved. This gives the precise argument for the last assertion of 2.17 and therefore its asphericity conclusion. The proof does not require $f$ to be fully faithful or cofinal.

The canonical factorization mentioned in the statement has the following form:
\[
 X\xrightarrow{p}Y/(f_!1_X)\xrightarrow{i}Y.
\]
Here $i$ is the discrete fibration associated with the presheaf, and
\[
 (f_!1_X)(y)=\pi _0(y\backslash X).
\]
The first map sends $x$ to the component represented by $(x,1_{f(x)})$. For an object $(y,c)$ of the category of elements, the comma category $(y,c)\backslash X$ is exactly the component $c$ of $y\backslash X$; it is nonempty and connected. Thus $p$ is cofinal in the terminology of these pages. The edition assigns ``\'etale'' to $p$, but it is the second map $i$ which is \'etale; the connected cofibres belong to $p$. The phrase ``objet $\cdots$ de $e_Y$'' similarly means an object of $\widehat Y$.

For the simplest example without a terminal object, take $X=Y=\mathbf N$ with its usual order and $f=1_X$. Its terminal presheaf is the increasing colimit of the representables, and its category of elements is again $\mathbf N$. This is filtered and weakly contractible, although no natural number is terminal. It supplies exactly the kind of example for which Grothendieck introduces the ind-object formulation. Corollary 2.18 is the same argument with an arbitrary presheaf in place of the terminal one.
\end{dcomment}


% Manuscript pp.36--53. Each block is independently anchored.
\begin{dcomment}{exposition}{p36:Mais je ne veux pas m'attarder à présent sur cette question.}{The quantifier in total asphericity}
Grothendieck is asking which indexing categories remain homotopically harmless after an ordinary change of index. The adjective ``total'' concerns every possible target of a functor; it does not merely strengthen the contractibility of the nerve of the original category. It is useful to isolate the precise condition before reading the proposed answer. Write $B A=|N A|$, and write $T(A)$ for the assertion that $B E$ is weakly contractible whenever $A\longrightarrow E$ is an ordinary cofinal functor. Here and below ``cofinal'' has the orientation fixed in Section~3: each category $e\backslash A$ must be nonempty and connected.

This formulation is equivalent to the one using $f_!1_A$ on the page. Indeed, for $f:A\longrightarrow Y$, put $P=f_!1_A$, so that
\[
 P(y)=\pi_0(y\backslash A).
\]
Let $Y/P$ be its category of elements. There is a canonical factorization
\[
 A\xrightarrow{q}Y/P\longrightarrow Y,
 \qquad a\longmapsto(fa,[a,1_{fa}]).
\]
For an object $(y,c)$ of $Y/P$, the comma category $(y,c)\backslash A$ is precisely the connected component $c$ of $y\backslash A$. Thus $q$ is cofinal. Conversely, if $f$ itself is cofinal, then $P=1_Y$ and $Y/P=Y$. Consequently $T(A)$ says exactly that the category of elements of every ordinary presheaf $f_!1_A$ is weakly contractible. The occurrence of $y\backslash Y$ in the printed question is a slip: it is $y\backslash A$ that must be connected.

A filtered category satisfies $T$. Its ordinary cofinal images are filtered by Corollary~3.4 below, and their nerves are weakly contractible. The difficult direction is whether this test forces the equalization of parallel arrows. The counterexample inserted on pages~36bis--36ter withdraws the attempted negative answer; it does not prove the converse. The converse holds for monoids and for categories with finitely many objects and fails in general; both are proved in the commentary ``Total asphericity is not filteredness'' below.

There is a terminological precaution for the later literature. In \citet[Definition~1.7.2 and Proposition~1.7.1]{Maltsiniotis2020}, a category is called totally aspherical when it is aspherical and its diagonal is aspherical, equivalently when products of representable presheaves are aspherical. This is a different quantifier from $T(A)$. In particular, the fact that the simplex category is a strict test category does not settle the question on this page. We retain $T(A)$ only inside this commentary, to keep the question itself unchanged.
\end{dcomment}

\begin{dcomment}{examples}{p36:p.ex. si on prend pour $X$ une catégorie-test}{The retracted proposition and a visible sphere}
The useful part of Proposition~2.19 is its elementary calculation with cofinal functors: the ordinary cofinal image of a nonempty sifted category is sifted. The unsupported step is the passage from that ordinary connectedness condition to weak contractibility. The edition marks this step and the proposition as false. The appeal in Counterexample~2.20 to Chapter~XII also needs its hypothesis restored: the reflection of a weak equivalence to the attaching category uses the strong form of $W(7\mathrm{bis})$, explicitly required in \citet[Chapter~XII, Section~V.3.4, Corollary~4, note~104]{Grothendieck1990XII}. That reflection is not valid for arbitrary weak homotopy equivalences. We can make the intended counterexample entirely explicit without using it.

The omitted reflection condition has a concrete obstruction. Choose a finite connected acyclic noncontractible simplicial complex $K$; a finite presentation complex with this homotopy type is given in \citet[Example~2.38 and Theorem~2C.5]{Hatcher2002}. Its suspension is simply connected and acyclic, hence contractible by the homology Whitehead theorem. The face posets of the two cones on $K$ are open subcategories covering the face poset of $\Sigma K$, with intersection the face poset of $K$. The two cones and their union are weakly contractible, but their intersection is not. Thus the weak equivalence from one cone into the union does not reflect to a weak equivalence from the intersection into the other cone. This contradicts the strong biconditional $W(7\mathrm{bis})$, while leaving the direct homotopy-pushout theorem intact.

Let $H=\Delta[1]/\partial\Delta[1]$, the simplicial circle, and form the category $E$ with full subcategory $\Delta$, one additional object $z$, and
\[
 E([n],z)=H_n,\qquad E(z,[n])=\varnothing,
 \qquad E(z,z)=\{1_z\}.
\]
Composition into $z$ is the simplicial action of $H$. Thus $\Delta/z=\Delta/H$, the category of simplices of $H$. Its classifying space has the homotopy type of $|H|=S^1$; this is the usual category-of-simplices realization, or the test-category equivalence for $\Delta$ \citep[Proposition~1.7.14]{Maltsiniotis2020}. Since $\Delta$ has terminal object $[0]$, its classifying space is contractible.

There is a direct description of $B E$. An $n$-simplex of $N E$ that does not lie in $N\Delta$ consists of a chain in $\Delta$ followed by an arrow into $z$, with possible repetitions of $z$. The arrow into $z$ and all its composites are exactly a simplex of $N(\Delta/H)$. These simplices glue to the mapping cone of the forgetful map $B(\Delta/H)\longrightarrow B\Delta$. Hence
\[
 B E\simeq\operatorname{hocolim}\bigl(B\Delta\longleftarrow
 B(\Delta/H)\longrightarrow *\bigr)\simeq\Sigma S^1\simeq S^2.
\]
The cone description can also be read degree by degree: adjoining the last vertex $z$ supplies the cone direction, and deleting that vertex applies the forgetful map. It therefore identifies the attaching maps, not merely the vertices of the two spaces.

The inclusion $i:\Delta\hookrightarrow E$ is ordinary initial: $\Delta/e$ has a terminal object when $e\in\Delta$, while $\Delta/z=\Delta/H$ is nonempty and connected. Thus $i^{\mathrm{op}}:\Delta^{\mathrm{op}}\longrightarrow E^{\mathrm{op}}$ is cofinal. The source $\Delta^{\mathrm{op}}$ is sifted: for $[m],[n]$, its category of common targets is opposite to $\Delta/(\Delta[m]\times\Delta[n])$, which is connected because the product of the two simplices is contractible. The target $E^{\mathrm{op}}$ is consequently sifted as well, whereas $B(E^{\mathrm{op}})\simeq S^2$. This supplies both a counterexample to the implication used in Proposition~2.19 and a direct counterexample to its asserted corollary about $\Delta^{\mathrm{op}}$.

The distinction is the same one encountered earlier in the chapter between products of connected objects and their higher homotopy. Here its failure has a specific witness, a two-dimensional sphere. It is not a failure of the composition or cancellation properties of ordinary cofinal functors.
\end{dcomment}

\begin{dcomment}{development}{p36:pourrait être valable.}{A necessary directedness condition}
Grothendieck asks, before returning to cofinality, whether his condition at least forces the existence of common successors. This weaker question admits an elementary answer. The argument below also settles his full question for preorders. We did not find this particular implication isolated in the consulted treatments of sifted categories, asphericity, or lax epimorphisms; the following is our continuation of the argument, with no assertion of priority.

\textbf{Proposition (the common-successor obstruction).}\ours
\textit{If $T(A)$ holds, then $A$ is nonempty and any two objects of $A$ admit arrows to a common object. Consequently, for a preorder $A$, the condition $T(A)$ is equivalent to filteredness.}

\textit{Proof.}
The identity of $A$ is cofinal, so $T(A)$ implies that $B A$ is weakly contractible. In particular $A$ is nonempty and connected. Suppose that objects $a,b$ have no common successor. Let $U$ be the full subcategory of objects receiving an arrow from $a$, and let $V$ be the full subcategory of objects receiving an arrow from $b$. These subcategories are disjoint and closed under passage to successors.

Let $P$ be the three-object category with precisely the nonidentity arrows
\[
 0\xrightarrow{p}1,\qquad 0\xrightarrow{q}2.
\]
Define $f:A\longrightarrow P$ by sending $U$ to $1$, $V$ to $2$, and all other objects to $0$. Closure under successors makes this a functor: an arrow cannot leave $U$ or $V$, and there are no arrows between them. The category $0\backslash A$ computed for $f$ is $A$ itself. The other two comma categories are $U$ and $V$. They are connected, since every object of $U$ receives an arrow from $a$, and every object of $V$ receives one from $b$. They need not have initial objects; the existence of those arrows is all that is used. Hence $f$ is cofinal.

We next give $P$ a cofinal image with nonzero second homology. Adjoin an object $z$ and arrows
\[
 r,s:z\rightrightarrows0,
 \qquad v:z\longrightarrow1,
 \qquad w:z\longrightarrow2,
\]
subject to $pr=ps=v$ and $qr=qs=w$, and no other nonidentity arrows. Call the resulting category $Q$. The inclusion $j:P\hookrightarrow Q$ is cofinal. For an object of $P$, the relevant comma category has an initial object. For $z$, the category $z\backslash P$ has four objects $r,s,v,w$ and four arrows
\[
 r\longrightarrow v,\quad r\longrightarrow w,
 \quad s\longrightarrow v,\quad s\longrightarrow w;
\]
it is nonempty and connected.

To see exactly what obstructs contractibility, use normalized simplicial chains of $NQ$. There are four nondegenerate two-simplices, denoted $(r,p),(s,p),(r,q),(s,q)$, and no nondegenerate simplices in higher dimensions. Their boundaries are
\[
\begin{aligned}
 \partial(r,p)&=p-v+r,&\partial(s,p)&=p-v+s,\\
 \partial(r,q)&=q-w+r,&\partial(s,q)&=q-w+s.
\end{aligned}
\]
Therefore
\[
 c=(r,p)-(s,p)-(r,q)+(s,q)
\]
is a nonzero cycle. Since there are no three-simplices, it is not a boundary. More precisely, the coefficients of $p,q,r,s,v,w$ in the boundary show that every two-cycle is an integral multiple of $c$, so $H_2(BQ;\mathbf Z)\cong\mathbf Z$. The composite $jf$ is cofinal, contradicting $T(A)$. This proves the common-successor assertion.

In a preorder, parallel arrows are already equal. Nonemptiness and common successors therefore imply filteredness. The converse follows from preservation of filteredness by ordinary cofinal images, as in Corollary~3.4. This proves the last assertion.\hfill$\square$

For general categories the argument leaves precisely the parallel-arrow equalization problem. The distinction is substantive: a one-object group category satisfies the common-successor condition for every group, while filteredness forces that group to be trivial. The next commentary settles it: the implication $T(A)\Rightarrow A\text{ filtered}$ holds when some object receives an arrow from every object, and fails for the category of finite ordinals and increasing injections.
\end{dcomment}

\begin{dcomment}{development}{p36:pourrait être valable.}{Total asphericity is not filteredness}
The question whether $T(A)$ forces $A$ to be filtered has a negative answer in general and a positive one for monoids and for categories with finitely many objects. Both rest on a description of $T(A)$ in which no target category appears. For a functor $H:A\to\mathrm{Set}$ write $\operatorname{El}(H)$ for its category of elements: the objects are the pairs $(a,h)$ with $h\in H(a)$, and an arrow $(a,h)\to(a',h')$ is an arrow $v:a\to a'$ with $H(v)h=h'$. The statements below were not found in the sources consulted, except where a reference is given; no priority is claimed.

\textbf{Theorem (a criterion on set-valued functors).}\ours
\textit{For a small category $A$ the following are equivalent: (i) $T(A)$; (ii) $A$ is nonempty and connected, and for every functor $H:A\to\mathrm{Set}$ every connected component of $\operatorname{El}(H)$ has the integral homology of a point; (iii) the same with ``is weakly contractible''. If they hold, every connected component $E$ of every $\operatorname{El}(H)$ satisfies $T(E)$.}

\textit{Proof.}
Let $K:A\to\mathrm{Set}$ have nonempty connected $\operatorname{El}(K)$, and let $A_K$ be $A$ with one object $x$ adjoined, $A_K(x,a)=K(a)$, no arrows into $x$ but its identity, and composition given by the action of $A$ on $K$. For the inclusion $A\to A_K$ the category $x\backslash A$ is $\operatorname{El}(K)$, and $a\backslash A$ has an initial object for $a\in A$; so the inclusion is cofinal. A chain of arrows of $A_K$ not contained in $A$ consists of $x$ (possibly repeated) followed by a chain of $\operatorname{El}(K)$. Hence $N(A_K)$ is $NA$ with the cone on $N\operatorname{El}(K)$ attached along the forgetful map, and
\[
 \widetilde H_{j+1}(BA_K)\longrightarrow\widetilde H_j(B\operatorname{El}(K))\longrightarrow\widetilde H_j(BA)
\]
is exact. Under (i) both $BA$ and $BA_K$ are weakly contractible, so $\operatorname{El}(K)$ has the homology of a point. A component of $\operatorname{El}(H)$ for an arbitrary $H$ is $\operatorname{El}(K)$ for the subfunctor $K$ of the elements lying in it. This proves (i)$\Rightarrow$(ii).

(ii)$\Rightarrow$(iii). Let $E$ be a component of $\operatorname{El}(H)$ and $L:E\to\mathrm{Set}$ a functor. Then $\operatorname{El}(L)\cong\operatorname{El}(H')$ for the functor $H'(a)=\coprod_hL(a,h)$ on $A$, the sum over the $h\in H(a)$ with $(a,h)\in E$. So $E$ again has property (ii). Let $\Pi E$ be the groupoid obtained by inverting all arrows of $E$, fix an object $e_0$, and identify $G=\operatorname{Aut}_{\Pi E}(e_0)$ with $\pi_1(BE,e_0)$. For a subgroup $\Gamma\subseteq G$ the functor $L_\Gamma(e)=\operatorname{Hom}_{\Pi E}(e_0,e)/\Gamma$ has bijective transition maps, and $\operatorname{El}(L_\Gamma)\to E$ is the connected covering with fundamental group $\Gamma$. By (ii) for $E$ its first homology $\Gamma^{\mathrm{ab}}$ vanishes. Taking $\Gamma$ cyclic gives $G=1$. A simply connected CW complex with the homology of a point is weakly contractible.

(iii)$\Rightarrow$(i). Let $f:A\to B$ be cofinal. For $b\in B$ the category $b\backslash A$ is $\operatorname{El}(a\mapsto B(b,f(a)))$; it is nonempty and connected, hence weakly contractible by (iii), and Quillen's Theorem~A makes $Bf$ a weak equivalence. The constant functor with value a point shows that $BA$ is weakly contractible. So is $BB$. The last assertion was shown in the second step.\hfill$\square$

In particular the integral-homology form of $T(A)$ (every cofinal image has the homology of a point) is equivalent to $T(A)$. With coefficients in a nonzero commutative ring $R$ the same proof, with the homology form of Theorem~A, gives: every cofinal image of $A$ has the $R$-homology of a point if and only if $A$ is nonempty and connected and the components of all $\operatorname{El}(H)$ do. For $R=\mathbf Q$ this is weaker: a finite group satisfies it.

\textbf{Proposition (exact colimits suffice).}\ours
\textit{If $A$ is nonempty and connected and $\varinjlim_A$ is exact on functors $A\to\mathrm{Ab}$, then $T(A)$.}

\textit{Proof.}
For $H:A\to\mathrm{Set}$ and $F:\operatorname{El}(H)\to\mathrm{Ab}$ put $(\Sigma_HF)(a)=\bigoplus_{h\in H(a)}F(a,h)$. The functor $\Sigma_H$ is exact and $\varinjlim_{\operatorname{El}(H)}F=\varinjlim_A\Sigma_HF$. So colimits over $\operatorname{El}(H)$ are exact, and so are colimits over a component $E$ (extend by zero). The left derived functors of $\varinjlim_E$ on the constant functor $\mathbf Z$ are the groups $H_j(BE)$; they vanish for $j>0$. Criterion (ii) applies.\hfill$\square$

\textbf{Theorem (a category with $T$ that is not filtered).}\ours
\textit{Let $\Delta'$ be the category of the ordered sets $[n]=\{0,\dots,n\}$, $n\geq0$, and strictly increasing maps. Then $T(\Delta')$ holds and $\Delta'$ is not filtered; it is not sifted either.}

\textit{Proof.}
The two maps $[0]\to[1]$ are not equalized by any injection, so $\Delta'$ is not filtered. A pair of maps $[0]\to[N]\leftarrow[0]$ is two points of $[N]$, and whether the first is smaller than, equal to or larger than the second is preserved by increasing injections; so the category of such pairs has three components and the diagonal of $\Delta'$ is not cofinal.

For exactness we show: \emph{for every $n$ there are $M$ and an integral combination $t=\sum_gt_g\,g$ of maps $g:[n]\to[M]$ of $\Delta'$, with $\sum_gt_g=1$, such that $t\circ h$ is the same combination for all maps $h:[m]\to[n]$ with the same $m$.} For $n=1$, $M=2$, one can take $t=(0,1)-(0,2)+(1,2)$, naming a map by its values. In general, a map $g:[n]\to[M]$ is determined by its gaps $e_0=g(0)$, $e_i=g(i)-g(i-1)-1$, $e_{n+1}=M-g(n)$, that is, by a monomial $z^e$ of degree $M-n$ in $z_0,\dots,z_{n+1}$; a combination $t$ becomes a homogeneous polynomial $P_t$, and $\sum t_g=P_t(1,\dots,1)$. Let $D_i$ be the substitution $z_{i+1}\mapsto z_i$ followed by renumbering to variables $y_0,\dots,y_n$. Omitting the value $g(i)$ merges two gaps and the omitted point, so $P_{t\delta^i}=y_iD_i(P_t)$ for the map $\delta^i:[n-1]\to[n]$ that omits $i$. We construct, for all $n\geq0$ and $k\geq1$, a homogeneous polynomial $C_{n,k}$ with integer coefficients, $C_{n,k}(1,\dots,1)=1$, such that $y_i^kD_i(C_{n,k})$ does not depend on $i$. Put $C_{0,k}=1$. Given $Q=C_{n-1,2k}$ in the variables $y_0,\dots,y_n$, we want $D_i(C_{n,k})=F_i:=Q\prod_{l\neq i}y_l^k$ for all $i$; then $y_i^kD_i(C_{n,k})=Q\prod_ly_l^k$. In the coordinates $x_i=z_i-z_{i+1}$ ($0\leq i\leq n$), $u=z_{n+1}$, the map $D_i$ is reduction modulo $x_i$. The prescriptions agree modulo $(x_i,x_j)$: for $j=i+1$ both give $D_i(Q)\prod_lw_l^k$, and for $j\geq i+2$ they give $w_{j-1}^{2k}D_{j-1}(Q)\prod_{l\neq i,j-1}w_l^k$ and $w_i^{2k}D_i(Q)\prod_{l\neq i,j-1}w_l^k$, which are equal by the property of $Q$. For a nonempty set $S$ of indices let $\Phi_S$ be the common value of the $F_i$, $i\in S$, after setting $x_i=0$ for all $i\in S$ (in homogeneous lifts written in the coordinates $x,u$). Then
\[
 C_{n,k}=\sum_{S\neq\varnothing}(-1)^{|S|+1}\Phi_S
\]
reduces to $F_j$ modulo $x_j$, because the terms $S$ and $S\cup\{j\}$, $j\notin S$, cancel modulo $x_j$ and the term $S=\{j\}$ remains. It is homogeneous, and its value at $(1,\dots,1)$, a point where all $x_i$ vanish, is $F_0(1,\dots,1)=1$. Now let $t$ correspond to $C_{n,1}$. Then $s=t\delta^i$ is independent of $i$, and from $\delta^{b+1}\delta^a=\delta^a\delta^b$ ($a\leq b$) we get $s\delta^a=s\delta^b$; by induction on $n$, $s\circ h'$ depends only on the source of $h'$, and every non-identity $h:[m]\to[n]$ factors through some $\delta^i$. This proves the claim.

Let $F\subseteq F'$ be functors $\Delta'\to\mathrm{Ab}$. Any two objects of $\Delta'$ map to a third, so every element of $\varinjlim F$ is the class of one element $\xi\in F([n])$. Suppose its image in $\varinjlim F'$ is zero. The colimit is the quotient of $\bigoplus_mF'([m])$ by the elements $F'(v)\eta-\eta$, so $\xi=\sum_k(F'(v_k)\eta_k-\eta_k)$ in this direct sum, for finitely many $v_k:[s_k]\to[r_k]$ and $\eta_k\in F'([s_k])$. Choose $N$ and maps $d_m:[m]\to[N]$ for the finitely many $m$ that occur; applying $F'(d_m)$ on each summand gives $F'(d_n)\xi=\sum_k\bigl(F'(d_{r_k}v_k)-F'(d_{s_k})\bigr)\eta_k$ in $F'([N])$. Choose $t$ for $N$ and write $F'(t)=\sum_gt_gF'(g)$. Since $t\circ d_{r_k}v_k=t\circ d_{s_k}$, applying $F'(t)$ gives $F'(\tau)\xi=0$ for $\tau=t\circ d_n$, a combination of maps $[n]\to[M]$ with coefficient sum $1$. As $F\subseteq F'$, also $F(\tau)\xi=0$. The class of $F(g)\xi$ in $\varinjlim F$ is the class of $\xi$ for every $g$, so the class of $F(\tau)\xi$ is that of $\xi$; it is zero. So $\varinjlim_{\Delta'}$ preserves monomorphisms and is exact; $\Delta'$ is nonempty and connected, and the proposition applies.\hfill$\square$

That colimits of abelian groups over this category are exact is attributed to Isbell and Mitchell in \citet[A.5]{AdamekKoubekReiterman1981}, together with the general criterion: $\varinjlim_A$ is exact on abelian groups if and only if each component of the category with the same objects and the integral combinations of arrows of coefficient sum $1$ as morphisms is filtered. The withdrawn corollary on this page claimed $T$ for $\Delta^{\mathrm{op}}$, which is false; the category $\Delta'$ is an example of what that corollary was after.

\textbf{Theorem (monoids; finitely many objects).}\ours
\textit{Let $A$ have an object $q$ that receives an arrow from every object. Then $T(A)$ holds if and only if $A$ is filtered. This applies to every monoid, and to every category with finitely many objects, where both $T(A)$ and filteredness provide such an object $q$.}

\textit{Proof.}
Assume $T(A)$. \emph{Step 1: for parallel arrows $a,b:p\to q'$ there are an object $r$ and an integral combination $t$ of arrows $q'\to r$, with coefficient sum $1$, such that $ta=tb$.} Let $P_c=\mathbf Z[A(c,-)]$, let $d:P_{q'}\to P_p$ be $c\mapsto ca-cb$, with kernel $K$ and image $J$, and let $H$ be the quotient of $A(p,-)$ by the relations $ca\sim cb$, so that $0\to J\to P_p\to\mathbf Z[H]\to0$ is exact. The functors $P_c$ are projective with $\varinjlim P_c=\mathbf Z$, and $\varinjlim d=0$. So $L_1\varinjlim\mathbf Z[H]$ is the kernel of $\varinjlim J\to\varinjlim P_p$; this map is zero, because $\varinjlim P_{q'}\to\varinjlim J$ is surjective and $\varinjlim d=0$. Hence $L_1\varinjlim\mathbf Z[H]=\varinjlim J=\operatorname{coker}(\varinjlim K\to\mathbf Z)$, the image of $\varinjlim K$ being generated by the coefficient sums of the elements of the groups $K(r)$. The left side is $H_1(B\operatorname{El}(H))$, and $\operatorname{El}(H)$ is connected (every object receives an arrow from the class of $\mathrm{id}_p$); by the criterion it vanishes. So $1$ is a sum of coefficient sums of finitely many elements of groups $K(r_i)$; pushing them to a common successor $r$ of the $r_i$ (the common-successor proposition of the preceding commentary) and adding gives $t$.

\emph{Step 2: the monoid $M=A(q,q)$ is filtered.} For $x,m\in M$, Step~1 for the pair $x,xm$, followed by an arrow $r\to q$, gives $t\in\mathbf Z[M]$ with coefficient sum $1$ and $tx=txm$. The element $\nu=tx$ of $\mathbf Z[Mx]$ is invariant under right multiplication $\rho$ by $m$. Let $S$ be its finite support. Each point of $S$ has a $\rho$-preimage in $S$, distinct points have distinct preimages, so $\rho$ permutes $S$ and $\nu$ is constant on its cycles; the coefficient sum $1$ is thus an integral combination of cycle lengths. Among all points of $Mx$ that are $\rho$-periodic choose one, $w$, of least period $\pi$. If $\pi>1$, let $\ell$ be a prime dividing $\pi$ and apply the same to $w$ in place of $x$: some cycle in $Mw$ has length prime to $\ell$, and its points $w'$ satisfy $w'm^\pi=w'$ too, so their period is a proper divisor of $\pi$, a contradiction. Hence there is $y$ with $yxm=yx$. Given $m,m'\in M$, take $y$ with $ym=y$ (case $x=1$), then $y'$ with $y'ym'=y'y$; then $y'ym=y'y=y'ym'$.

\emph{Step 3.} Let $a,b:p\to q'$ be parallel, choose $g:q'\to q$, and by Step~1 an element $t=\sum n_ic_i\in\mathbf Z[M]$ of coefficient sum $1$ with $tga=tgb$. In the filtered monoid $M$ there is $y$ with $yc_i=y$ for all $i$. Then $yga=\sum n_iyc_iga=\sum n_iyc_igb=ygb$. With common successors this shows that $A$ is filtered. The converse is the fact that filtered categories satisfy $T$.\hfill$\square$

So Grothendieck's expectation is right in every finite case and for all monoids, and the smallest counterexamples have infinitely many objects. For the monoid $\{1,e,f\}$ with $uv=v$ for $u,v\in\{e,f\}$, which is cofiltered with contractible classifying space, the theorem says that $T$ fails; a cofinal image of the homotopy type of $S^2$ is obtained by adjoining a source $x$ with two arrows $h,i:x\to*$ on which $e$ and $f$ act by the constant map with value $i$. Whether $T(A)$ implies exactness of $\varinjlim_A$ on abelian groups, the converse of the proposition, is not decided here.
\end{dcomment}

\begin{dcomment}{exposition}{p37}{The represented derivator at dimension zero}
Grothendieck now uses sets to separate the genuinely categorical part of cofinality from its higher homotopical refinements. The equivalence $\HOT_{W_0}(I)\simeq\operatorname{Fun}(I^{\mathrm{op}},\mathrm{Set})$ has a short proof that explains why no hidden higher information survives. Let $d:\mathrm{Set}\longrightarrow\mathrm{Cat}$ assign the discrete category, and let $\pi_0$ be connected components. There is an adjunction $\pi_0\dashv d$, with $\pi_0d=1$. For each category $C$, the unit $C\longrightarrow d\pi_0C$ is a $W_0$-equivalence. Applied pointwise to an $I^{\mathrm{op}}$-diagram, these units give a natural objectwise $W_0$-equivalence from that diagram to its discrete component diagram. Thus the two functors become inverse after localization. Since the construction commutes with restriction, it gives the stated equivalence of diagram theories.

The description by $\mathrm{Cat}/I$ requires the relative, rather than the global, component construction. To $X\longrightarrow I$ one associates the presheaf $i\longmapsto\pi_0(i\backslash X)$. Localizing the slice merely at functors that induce a bijection on the components of their total categories would give a different category. The relevant class is detected in all these comma categories, exactly as in the relative localization recalled in the commentary on Chapter~I, manuscript p.~43 (``What the localization axioms do and do not say'').

The large table then has an economical interpretation. The words ``aspherical'' and ``coaspherical'' become ordinary initiality and finality; the proper and smooth conditions say that specified strict lift categories are nonempty and connected. Right and left Kan extensions retain their ordinary meanings. Nothing in this reduction licenses replacing connected lift categories by contractible ones. Conversely, it makes the degree-zero assertions accessible by elementary calculations with connected components, without invoking a construction of general homotopy Kan extensions.
\end{dcomment}

\begin{dcomment}{exposition}{p40}{The lift categories behind ordinary base change}
The base-change assertion is an application of the lift-category criterion already developed in the commentary on Chapter~I, manuscript pp.~68--70 (``The category of lifts behind properness'' and ``Completing the tests in the proper base-change theorem''). Its present specialization is useful because all the relevant tests are tests of connectedness. For $u:A\longrightarrow B$, properness means that, for every $a\in A$ and arrow $\beta:u(a)\longrightarrow b$, the category
\[
 L(a,\beta)=\{a\xrightarrow{h}a'\mid u(a')=b, u(h)=\beta\}
\]
of strict lifts, with arrows over $1_b$, is nonempty and connected. Smoothness is the dual condition on arrows into $u(a)$. An opfibration supplies initial objects in the displayed lift categories; a fibration supplies the dual terminal objects. Connectedness is the exact weakening that suffices for ordinary Kan extensions.

For a strict pullback square as in (3.9), the comparison of the comma categories computing a right Kan extension is final if $u$ is proper or if the horizontal base-change functor is smooth. The ordinary cofinality theorem therefore gives the right base-change isomorphism. Reversing arrows gives the left formula. The same argument works in any coefficient category in which the requisite comma limits exist, so that these Kan extensions are computed pointwise. One can apply $\operatorname{Hom}(m,-)$ to the limits, or dually $\operatorname{Hom}(-,m)$ to the colimits, and reduce to the set-valued cofinality theorem. This avoids imposing unnecessary completeness assumptions on the coefficient category.

The passage from a $W_0$-equivalence to an equivalence measured by constant coefficients has a similarly elementary explanation. A limit of a constant diagram is a product indexed by the connected components of its shape, and a colimit is the corresponding coproduct, whenever those objects exist. A bijection on components preserves both. The printed footnote correctly excludes the two fibration notions from the blanket transfer of terminology: preserving a larger class of equivalences under every base change is a stronger requirement, even when the equivalences themselves have a simpler definition. The next proposition measures this effect exactly.
\end{dcomment}

\begin{dcomment}{exposition}{p42:Cette condition est aussi suffisante}{Why arrows detect universal connectedness}
Proposition~3.1 says that universal preservation of connected components can be checked on vertices and single arrows. Its proof is a useful model for reading the much stronger componentwise condition in part~(b). We give the details because the proposition is stated, rather than proved, here.

Suppose first that $u:A\longrightarrow B$ is a universal $W_0$-equivalence. Pullback along any $[1]\longrightarrow B$ gives a connected nonempty category, since $[1]$ is connected. Conversely, assume this arrow test. Applying it to an identity arrow at $b$ shows that $A_b\times[1]$, and hence $A_b$, is connected and nonempty. Let $v:J\longrightarrow B$ be arbitrary. Every object of $J$ has a nonempty fibre in $A\times_BJ$, so the map on components is surjective. If two objects of the pullback project to objects joined by a zigzag in $J$, each edge of that zigzag has a connected pullback by hypothesis. Its endpoints can be joined through that pullback, and the connected vertex fibres join any two choices of endpoints. Concatenating these finite zigzags proves injectivity. Thus
\[
 \pi_0(A\times_BJ)\xrightarrow{\sim}\pi_0(J)
\]
for every $J$, which is the required universal equivalence. Notice that the test concerns strict pullbacks, not comma categories.

Now let $A_\lambda$ be a connected component of $A$, mapping to the component $B_\mu$. If pullback by $u$ preserves every $W_0$-equivalence, apply it to $J\amalg\coprod_{\nu\ne\mu}B_\nu\longrightarrow B$, where $J\longrightarrow B_\mu$ is any connected nonempty category over $B_\mu$. The resulting map on components forces $A_\lambda\times_{B_\mu}J$ to be connected and nonempty. Decomposing an arbitrary category over $B_\mu$ into its components then shows that $A_\lambda\longrightarrow B_\mu$ is a universal $W_0$-equivalence. Conversely, if all these component maps are universal equivalences, then for any $J\longrightarrow B$ there is a natural bijection
\[
 \pi_0(A\times_BJ)\cong
 \pi_0(A)\times_{\pi_0(B)}\pi_0(J).
\]
This formula immediately shows that pullback by $u$ preserves every bijection on connected components. It proves part~(b), including components of $B$ not met by $A$.

The necessary conditions in part~(c) come from the same tests with categorical homotopy equivalences. The inclusion of the initial object $1_b$ into $b\backslash B$ is such an equivalence; its pullback identifies $A_b\longrightarrow b\backslash A$ on components. Testing both endpoints of an arrow in $B$ makes the resulting component presheaf locally constant. When $u$ is smooth, ordinary left base change identifies the component presheaf after every base change with the restriction of this local system. Its colimit depends only on the fundamental groupoid of the indexing category, and categorical homotopy equivalences induce equivalences of those groupoids. This proves the stated sufficiency. The opposite-category assertion follows directly because categorical homotopies and $\pi_0$ are unchanged by reversal of arrows.
\end{dcomment}

\begin{dcomment}{examples}{p43:i.e. que $u$ soit un isomorphisme}{Monodromy explains the strictness}
The covering example isolates what the componentwise test excludes. This is the same monodromy issue as in the commentary on Chapter~I, manuscript p.~98 (``Sets already detect the missing monodromy''), now applied to the two precise definitions in (3.8). Let $G$ be a nontrivial group. Its translation groupoid $E G$ has object set $G$ and exactly one arrow between each ordered pair of objects. The functor
\[
 E G\longrightarrow B G
\]
sends $g\longrightarrow h$ to $hg^{-1}$. It is a discrete fibration and a discrete opfibration, so the strict lift categories are singletons and the map is proper and smooth. Both total and base categories are connected, but the strict fibre is the discrete set $G$.

The inclusion of the sole object $*\longrightarrow BG$ is a $W_0$-equivalence. Its pullback along $EG\longrightarrow BG$ is the map $G_{\mathrm{disc}}\longrightarrow EG$, whose map on components is $G\longrightarrow *$. This proves directly that the covering is not a $W_0$-fibration. It is a weak $W_0$-fibration by the smooth criterion of Proposition~3.1(c): the relative component presheaf is the locally constant $G$-set with regular action. Thus weak fibrations permit monodromy between different fibre components, whereas the strong $W_0$ condition treats each connected component of the total category separately.

There is a subscript caution in the preceding printed paragraph. The conclusion established by this example is ``not necessarily a $W_0$-fibration.'' It does not assert that coverings fail to preserve weak homotopy equivalences under pullback. In fact their nerves are simplicial coverings, and they do preserve weak homotopy equivalences. Similarly a $W_0$-perfect map gives a weak $W_0$-fibration; a $W$-perfect map gives the corresponding weak $W$-fibration. These are distinct assertions when $W\subsetneq W_0$.
\end{dcomment}

\begin{dcomment}{exposition}{p45:i.e. que $i\bs I'$ est 0-connexe, q.e.d.}{What the ordinary cofinality theorem tests}
Grothendieck has now proved the criterion announced on the opening page, with the sufficient directedness condition replaced by the exact connectedness condition. The proof is the representability calculation already explained in ``The presheaf represented by a change of index'': the functor $F\longmapsto\Gamma(I',u^*F)$ is represented by $u_!1_{I'}$, and
\[
 (u_!1_{I'})(i)=\pi_0(i\backslash I').
\]
Consequently the comparison on all set-valued presheaves is an isomorphism precisely when $u_!1_{I'}\longrightarrow1_I$ is an isomorphism. Applying ordinary Hom functors reduces arbitrary existing limits and colimits to this case. This is the classical theorem of SGA~4, Expos\'e~I, Section~8.1, in the variance used throughout these pages \citep{SGA4I1972}.

Two small examples distinguish its quantifiers. The inclusion $\{1\}\hookrightarrow[1]$ is cofinal, because each $i\backslash\{1\}$ is a point. For a presheaf $F(1)\longrightarrow F(0)$, its limit is therefore $F(1)$. But the restriction functor forgets $F(0)$ and is not fully faithful. On the other hand, for any nonempty connected category $C$, the projection $C\longrightarrow *$ is cofinal even when $BC$ has nontrivial homotopy groups. Ordinary cofinality is therefore stronger than a bijection on components of arbitrary source and target, weaker than full faithfulness of restriction, and weaker than homotopy cofinality in a different direction. These comparisons refer to different properties of a functor, not to alternative conventions for the same word.

Footnote~53 distinguishes two degree-zero conditions on a morphism of topoi $f:\mathcal E'\longrightarrow\mathcal E$. Its coasphericity condition preserves global sections of every object. Its asphericity condition instead requires $f^*U$ to be noninitial and connected whenever $U$ is. For locally connected topoi, the latter is equivalent to preservation of component sets: decompose $U$ into its noninitial connected components and use preservation of coproducts by $f^*$. For the morphism induced by $u:I'\longrightarrow I$, the restricted representable $u^*h_i$ has category of elements $I'/i$. Thus this asphericity condition is ordinary initiality of $u$, by the cofinality theorem with arrows reversed; the section condition is finality. The two endpoints of $[1]$ distinguish them: inclusion of $0$ is initial, while inclusion of $1$ is final.

The supplementary Yoneda test in footnote~54 should use the functor $I\longrightarrow I^{\wedge}$. Its colimit is $1_I$, whereas its restriction along $u$ has colimit $u_!1_{I'}$. The printed $I'\longrightarrow I'^{\wedge}$ has the wrong domain for the test as stated. The further suggestion in footnote~53 about higher cohomology should be read with coefficients specified: preservation of all set-valued sections, preservation of abelian cohomology through a fixed degree, and equivalence on all homotopy-valued diagrams detect different truncations. The earlier cohomological cofinality block supplies the derived comparison; the ordinary theorem here establishes its degree-zero model.
\end{dcomment}

\begin{dcomment}{development}{p49:La vérification de d) par AQT}{The filtered criterion and eventual lifting}
Corollary~3.4 converts connectedness of comma categories into the two concrete operations needed for filtered colimits: reaching the image of $u$, and equalizing parallel arrows after reaching it. The proof of parts~(b) and~(c) is present on the page. There is one harmless slip in the statement of~(c): its last assertion concerns filteredness of $I$, not the already assumed filteredness of $I'$. Part~(d) is left under the abbreviation AQT. Its proof uses full faithfulness in two different places.

One variance qualification is essential. Definition~2.2 called the completion of a cospan by a common predecessor PS~1 and explicitly called its right-hand dual ``colocalement faiblement filtrante.'' Here part~(d) and footnote~57 instead use ``pr\'efiltrante'' or ``localement faiblement filtrante'' for the right-hand condition: a span must admit a commuting square with a common successor. The proof requires that interpretation. Under the earlier left-hand reading, part~(d) would be false: the poset $I=(0\longrightarrow1,\,0\longrightarrow2)$ satisfies left PS~1, and the full inclusion $\{1,2\}\hookrightarrow I$ satisfies F~1, but its comma category under $0$ is a discrete two-point category. Thus this is a necessary clarification of the printed terminology, not a variance-neutral use of PS~1.

Assume that $I$ satisfies the right-hand Ore condition PS~1, that $u:I'\longrightarrow I$ is fully faithful, and that F~1 holds. Given objects $(a,\alpha:i\longrightarrow u a)$ and $(b,\beta:i\longrightarrow u b)$ of $i\backslash I'$, PS~1 supplies arrows $r:u a\longrightarrow j$ and $s:u b\longrightarrow j$ with $r\alpha=s\beta$. Choose $j\longrightarrow u c$ by F~1. Fullness lifts the two resulting arrows $u a\longrightarrow u c$ and $u b\longrightarrow u c$ to arrows in $I'$, and they give a common successor of the original two objects in $i\backslash I'$. F~1 also makes that comma category nonempty. It is therefore weakly filtered, hence connected, and $u$ is cofinal.

To transfer PS~1 itself to $I'$, start with a span $a\longrightarrow b$, $a\longrightarrow c$ in $I'$. Complete its image in $I$ to a commuting square, reach an object in the image of $u$ by F~1, and lift the two new arrows by fullness. Faithfulness ensures that the lifted square commutes. Thus $I'$ satisfies PS~1. If $I$ is filtered, a parallel pair $r,s:a\rightrightarrows b$ in $I'$ can likewise be equalized in $I$, followed by an arrow to $u c$ and lifted to an equalizer arrow $b\longrightarrow c$ in $I'$. Faithfulness again supplies the equality. Nonemptiness follows from F~1, so $I'$ is filtered. The same reasoning, or the discrete-opfibration projection $i\backslash I'\longrightarrow I'$, makes each comma category filtered.

For locally filtered categories, the additional condition F~0 in footnotes~57--58 ensures that objects whose images lie in the same component already lie in the same source component. The arguments can then be carried out in that component. Without this component condition, F~1 and F~2 do not prevent several disconnected copies of one indexing category from being mapped onto it. This is why the global filtered statement and its local variant have different hypotheses.

Footnote~55 asks for a statement encompassing this corollary and Chapter~XVIII. The exact later result is Lemma~4.9.6.3, part~1, manuscript p.~224, whose proof is in \S4.9.6.7. Its additional condition is eventual lifting:
\[
 \text{F~3:}\quad
 \theta:u(a)\longrightarrow u(b)
 \ \Longrightarrow\
 \exists\,v:b\longrightarrow c,\ w:a\longrightarrow c,
 \qquad u(v)\theta=u(w).
\]
For a faithful functor into a filtered category, the lemma identifies F~1 and F~3 with filteredness of the source together with cofinality \citep[Lemma~4.9.6.3(1), pp.~81--82]{Grothendieck1990XVIII}. The following finite formulation combines the two arguments, retaining only the right Ore hypothesis where that suffices. It is our proof-level synthesis of the cited passages, without a claim of priority.

\textbf{A finite comparison with eventual lifting.}\ours
\textit{If the target satisfies right PS~1, F~1 and F~3 make $u$ cofinal. If $u$ is faithful, its source also satisfies right PS~1 and is filtered whenever the target is filtered.}

Indeed, given $(a,\alpha)$ and $(b,\beta)$ under $i$, right PS~1 followed by F~1 produces $\rho_a:u(a)\to u(c)$ and $\rho_b:u(b)\to u(c)$ with $\rho_a\alpha=\rho_b\beta$. Apply F~3 first to $\rho_a$, obtaining $v:c\to d$, $w:a\to d$, and then to $u(v)\rho_b$, obtaining $t:d\to e$, $z:b\to e$. The equality
\[
 u(tw)\alpha=u(t)u(v)\rho_a\alpha
              =u(t)u(v)\rho_b\beta=u(z)\beta
\]
gives a common successor in the comma category. F~1 gives nonemptiness, hence cofinality. For a span in the source, this construction makes its two composites equal after applying $u$; faithfulness makes them equal before applying it. If the target is filtered, equalize the images of any parallel pair, reach another image object by F~1, and lift the resulting arrow eventually by F~3. Faithfulness then gives an equalizing arrow in the source. Cofinality supplies connectedness and nonemptiness, completing filteredness. Conversely, when the source is filtered and $u$ cofinal, the already proved filteredness of $u(a)\backslash I'$ gives F~3 by taking a common successor of $(a,1)$ and $(b,\theta)$.

Full faithfulness makes F~3 automatic, recovering part~(d). The source-filtered case is the F~1--F~2 criterion in part~(c). Thus the finite comparison requested in footnote~55 is explicit. Part~2 of the later lemma strengthens F~3 to simultaneous lifting for families of bounded cardinality and imposes the corresponding larger filtering hypotheses; that cardinal refinement is a separate assertion, not a consequence of the finite conditions alone.
\end{dcomment}

\begin{dcomment}{exposition}{p51}{The missing pointwise limit hypothesis}
Grothendieck isolates a real difference between testing every coefficient on a slice and testing only coefficients restricted from the ambient category. There is also a second issue here: full faithfulness for coefficients in $\mathcal M$ does not imply full faithfulness after replacing $\mathcal M$ by its presheaf completion. The sufficient implication (4.8) is valid. The converse discussed on page~51, including its asserted ordered case, requires hypotheses ensuring that the relevant pointwise Kan-extension limits exist in $\mathcal M$.

Take $I=[1]=(0\longrightarrow1)$, $I'=\{1\}$, and let $\mathcal M$ be the discrete category on two objects $m,n$. Every $\mathcal M$-valued presheaf on $I$ is constant, so restriction $\mathcal M(I)\longrightarrow\mathcal M(I')$ is an equivalence. Nevertheless, at $i=0$ the induced functor on slices is $\varnothing\longrightarrow *$. In $\mathcal M^{\wedge}\simeq\mathrm{Set}\times\mathrm{Set}$, the Yoneda image of $m$ is $(1,\varnothing)$, whereas the limit of the empty diagram is $(1,1)$. Thus the comparison
\[
 \varprojlim_{I/0}y(m)=(1,\varnothing)
 \longrightarrow\varprojlim_{I'/0}y(m)=(1,1)
\]
is not invertible. This coefficient on $I/0$ even extends to $I$. Hence the intermediate assertion that $\mathcal M$-connectedness is equivalent to the slice comparison for restricted coefficients fails without pointwise limit hypotheses.

Existence of an abstract adjoint alone is insufficient: in this example restriction is an equivalence and has both adjoints. What fails is the formula that would compute its right adjoint by an empty limit inside $\mathcal M$. Completeness of $\mathcal M$, or existence of all the particular comma limits involved, supplies the needed repair.

The extension observation for posets is independently correct. Identify $I/i$ with $A=\{x\mid x\leq i\}$. The map $r:I\longrightarrow A$ sending $x$ to itself when $x\leq i$ and to $i$ otherwise is order preserving: no arrow can enter $A$ from outside, and an arrow leaving $A$ is sent to $x\leq i$. With $j:A\hookrightarrow I$, one has $rj=1_A$, so every coefficient on $A$ extends by $r^{\mathrm{op}}$. Under the pointwise limit hypotheses this proves the ordered-case converse. It does not remove those hypotheses. Proposition~4.1 subsequently quantifies over all coefficient categories and can use sets from the start; its intended theorem is unaffected by this coefficient-specific error.
\end{dcomment}

\begin{dcomment}{later}{p53}{Biasphericity, lax epimorphisms, and factorization categories}
The new question in Section~4 is whether a diagram and all its natural transformations can be recovered after restriction. The answer is governed by factorizations of arrows, just as ordinary cofinality was governed by objects of comma categories. This is the ordinary version of the factorization test encountered in the commentary on Chapter~I, manuscript p.~49 (``The factorization category controls base change''). The proof of Proposition~4.1 continues beyond these pages; we give the complete argument for its corrected statement here.

For $u:A\longrightarrow B$ and $\gamma:b_0\longrightarrow b_1$, let $\mathcal F_u(\gamma)$ have objects
\[
 (a,p,q),\qquad b_0\xrightarrow{p}u a\xrightarrow{q}b_1,
 \quad qp=\gamma,
\]
and arrows $h:a\longrightarrow a'$ satisfying $u(h)p=p'$ and $q'u(h)=q$. Then the following conditions are equivalent: restriction on presheaves is fully faithful for every coefficient category; it is fully faithful for sets; and every $\mathcal F_u(\gamma)$ is nonempty and connected.

To prove this, first work with sets, where both Kan extensions exist. Full faithfulness of $u^*$ is equivalent to the counit $u_!u^*F\longrightarrow F$ being an isomorphism. Both sides preserve colimits in $F$, and every presheaf is a colimit of representables, so it suffices to take $F=h_{b_1}=B(-,b_1)$. At $b_0$ the pointwise colimit formula gives
\[
 (u_!u^*h_{b_1})(b_0)
 \cong\coprod_{\gamma\in B(b_0,b_1)}
             \pi_0\mathcal F_u(\gamma).
\]
Indeed, the generators are pairs $b_0\longrightarrow u a\longrightarrow b_1$, and the colimit relations are exactly the zigzags induced by arrows of $\mathcal F_u(\gamma)$. Composition labels every such relation by the same $\gamma$. The counit sends a connected component to that composite. It is bijective if and only if each indicated category is nonempty and connected.

The set-valued criterion implies the version for arbitrary coefficients. It first implies full faithfulness for coefficients in every presheaf category, by applying the set-valued statement in each variable. A small category $\mathcal M$ embeds fully faithfully into its presheaf category by Yoneda; the full-faithfulness property descends to this full subcategory. For a possibly large locally small $\mathcal M$, any two diagrams on the small category $B$ take values in a small full subcategory. Applying the preceding argument there descends every natural transformation and proves uniqueness. This establishes the reductions (4.3)--(4.6), with the usual universe enlargement when necessary.

For fixed $b_1$, the comma category of the induced functor $A/b_1\longrightarrow B/b_1$ under $(b_0,\gamma)$ is $\mathcal F_u(\gamma)$. Thus condition~(d) is equivalent to the factorization criterion. Its dual says that
\[
 b_0\backslash A\longrightarrow b_0\backslash B
 \quad\text{is \emph{co-cofinal}, that is, initial.}
\]
The edition instead prints ``cofinal'' in~(d$'$). That cannot be retained as a mathematical equivalence. The inclusion $u:\{1\}\hookrightarrow[1]$ satisfies the printed condition: for $i=0$ it induces the final-vertex inclusion into $[1]$, and for $i=1$ it induces the identity of a point. Yet restriction on presheaves is not faithful: take $F(1)=\varnothing$, $F(0)=\{0,1\}$; its identity and the automorphism transposing $0,1$ become the same map after restriction. The corrected dual condition has comma categories opposite to $\mathcal F_u(\gamma)$, so it gives exactly the required connectedness. The two adjunction formulations and self-duality now follow, proving all the intended equivalences.

In later terminology these functors are \emph{lax epimorphisms}, also called absolutely dense functors. \citet{AdamekElBashirSobralVelebil2001} identify them by full faithfulness of set-valued restriction and by connected categories of factorizations; the latter characterization is restated explicitly in \citet[Theorem~4.1]{NunesSousa2022}. For derivators, \citet[Theorem~6.10]{GrothStovicek2016Trees} replace connectedness by weak contractibility of the same factorization categories and obtain \emph{homotopical epimorphisms}. Thus the arrowwise obstruction persists; the homotopy theory changes the required strength of its vanishing.

A quotient group displays the difference concretely. For a surjective homomorphism $\rho:H\twoheadrightarrow G$, restriction of $G$-sets to $H$-sets is fully faithful, because equivariance under $H$ is equivalent to equivariance under its image $G$. For any $g\in G$, the category $\mathcal F_{B\rho}(g)$ is the transitive action groupoid with object set $G$, on which $H$ acts through $\rho$; its stabilizer is $\ker\rho$. It is therefore equivalent to $B\ker\rho$. The functor $B\rho$ is an ordinary lax epimorphism, but it is a homotopical epimorphism only when $\ker\rho$ is trivial. With $\rho:\mathbf Z\twoheadrightarrow\mathbf Z/n\mathbf Z$, the obstruction has the homotopy type of a circle. The full-faithfulness statement about ordinary diagrams is correct; it simply does not recover that circle.
\end{dcomment}
